\documentclass[11pt]{amsart}
\usepackage[margin=1.1in]{geometry}
\usepackage{amsmath,amssymb,amsthm}
\usepackage{booktabs}
\usepackage{xcolor}
\usepackage[colorlinks=true,linkcolor=blue!50!black,citecolor=blue!50!black]{hyperref}

% Referee version: set to true to print line numbers (journal referee copy).
% Leave false for the arXiv/archival copy.
\newif\ifreferee
\refereefalse
\usepackage{lineno}
\ifreferee\linenumbers\fi

\newtheorem{theorem}{Theorem}[section]
\newtheorem{lemma}[theorem]{Lemma}
\newtheorem{proposition}[theorem]{Proposition}
\newtheorem{corollary}[theorem]{Corollary}
\theoremstyle{definition}
\newtheorem{definition}[theorem]{Definition}
\theoremstyle{remark}
\newtheorem{remark}[theorem]{Remark}
\newtheorem{convention}[theorem]{Convention}

\DeclareMathOperator{\Inn}{Inn}
\DeclareMathOperator{\Aut}{Aut}
\DeclareMathOperator{\Out}{Out}
\DeclareMathOperator{\Syl}{Syl}
\DeclareMathOperator{\Fix}{Fix}
\DeclareMathOperator{\PSL}{PSL}
\DeclareMathOperator{\PSU}{PSU}
\DeclareMathOperator{\PSp}{PSp}
\DeclareMathOperator{\Sp}{Sp}
\DeclareMathOperator{\SL}{SL}
\DeclareMathOperator{\SU}{SU}
\DeclareMathOperator{\GL}{GL}
\DeclareMathOperator{\GU}{GU}
\DeclareMathOperator{\Sz}{Sz}
\DeclareMathOperator{\ord}{ord}
\newcommand{\Z}{\mathbb{Z}}
\newcommand{\F}{\mathbb{F}}
\newcommand{\Pos}{\mathcal{P}}
\newcommand{\vp}[1]{v_{#1}}

\title[Groups factorized by non-conjugate maximal subgroups]{Finite groups
that are the product of every pair of non-conjugate maximal subgroups are
soluble}

\author{Richie Sater}
\address{Independent Researcher, United States}
\email{richiesater@gmail.com}
\urladdr{https://orcid.org/0009-0007-9051-8207}

\date{August 5, 2026}

\subjclass[2020]{Primary 20D40; Secondary 20E28, 20D05}
\keywords{Factorizations of finite groups, maximal subgroups, soluble
groups, simple groups, Kourovka Notebook Problem 10.34}

\begin{document}

\begin{abstract}
We prove that every finite group that is the product of every pair of
its non-conjugate maximal subgroups is soluble, answering Problem
10.34 of the Kourovka Notebook (V.~S.~Monakhov, 1986) in the negative.
Tikhonenko and Tyutyanov settled the almost simple case in 2010. What remains
is the possibility of a minimal counterexample whose unique minimal normal
subgroup is $S^k$, with $S$ non-abelian simple and $k\geq2$. Because $k$ is
unbounded, no bound depending on a fixed ratio of subgroup orders can handle
all such groups. We rule out this possibility with a divisibility criterion:
a single pair of automorphism-stable conjugacy classes of subgroups of $S$,
satisfying a valuation inequality, excludes the configuration for every
$k\geq2$ and every admissible embedding. Such pairs are constructed uniformly
for every infinite family of finite simple groups, with Zsigmondy primes as
the arithmetic obstruction. The sporadic groups and the finitely many
exceptional cases are handled by independently checkable GAP certificates.
\end{abstract}

\maketitle

\section{Introduction}\label{sec:intro}

Problem 10.34 of the Kourovka Notebook, posed by V.~S.~Monakhov in 1986
\cite{Kourovka}, asks:

\begin{quote}
\emph{Does there exist a finite non-soluble group that coincides with the
product of any two of its non-conjugate maximal subgroups?}
\end{quote}

Many soluble groups have this property: it holds vacuously whenever all
maximal subgroups are conjugate, and $S_4=AB$ for every pair of non-conjugate
maximal subgroups $A,B$. The content of the problem is whether the property
forces solubility. We prove that it does.

\begin{theorem}[Main Theorem]\label{thm:main}
Let $G$ be a finite group such that $G=AB$ for every pair of maximal
subgroups $A,B\leq G$ that are not conjugate in $G$. Then $G$ is soluble.
\end{theorem}

Throughout, we say $G$ has \emph{property $\mathrm{P}$} if $G=AB$ for
every pair of non-conjugate maximal subgroups $A,B\leq G$.

The almost simple case of Theorem~\ref{thm:main} was proved by Tikhonenko
and Tyutyanov \cite{TT2010}. We use their theorem as the essential $k=1$
branch of the main proof. The new content is the case $k\geq2$.
Tikhonenko and Tyutyanov also record that Zenkov had announced a negative
answer to the full problem in a 1997 conference abstract \cite{Zenkov1997},
while stating that no proof of that announcement was known to
them. Searches completed on 2 August 2026 did not identify a published proof,
and Problem 10.34 remains listed among the unsolved problems in the July 2026
revision of the Kourovka Notebook \cite{Kourovka}. We therefore claim a proof
of the theorem, not priority for the conclusion. For comparison,
Theorem~\ref{thm:Dprime} shows that the present criterion also recovers the
almost simple conclusion; it is not used in the proof of
Theorem~\ref{thm:main}.

The maximal factorizations of finite simple groups and their automorphism
groups were classified by Liebeck, Praeger and Saxl \cite{LPS1990}. Their
tables classify the factorizations that \emph{do} occur. Our route needs only
to exhibit, for each coordinate closure, one pair of classes whose associated
maximal subgroups cannot factorize $G$; the valuation obstruction avoids
importing the full maximal-factorization tables.

\subsection*{Strategy}
A minimal counterexample $G$ has a unique minimal normal subgroup
$N=S^k$ with $S$ non-abelian simple, and embeds in $\Aut(S)\wr S_k$
(\S\ref{sec:reduction}). For an automorphism-stable conjugacy class $[V]$
of suitable subgroups $V<S$, the normalizer $B_V=N_G(V^k)$ is a maximal
subgroup of $G$ of order exactly $|G/N|\cdot|V|^k$
(\S\ref{sec:machinery}). Property $\mathrm{P}$ forces $G=B_VB_W$ for
distinct such classes, and the resulting integrality condition
$|B_V\cap B_W|\in\Z$ produces, for a well-chosen prime $p$, the
divisibility $p^{\eta k}\mid |G/N|$ with $\eta$ independent of $k$, which is
incompatible with $|G/N|$ dividing $x^k\,k!$, where $x=|X/\Inn(S)|$ for
the relevant $X\leq\Aut(S)$ (\S\ref{sec:criterion}). The conclusion
(Theorem~\ref{thm:D}) is that a \emph{single} pair of stable classes in
$S$, satisfying a valuation inequality, rules out the corresponding
minimal-normal-subgroup configuration for every $k\geq 2$ and every
admissible embedding.

The $k\geq2$ branch thus reduces to exhibiting such a pair for every finite
simple group $S$. By the classification of finite simple groups
\cite[Ch.~1, Table~I, pp.~8--10]{GLS1} (see also
\cite[Thm.~1.1]{RoneyDougal2021}), the possibilities for $S$ comprise
finitely many infinite families, the $26$ sporadic groups, and the Tits group.
For the infinite families the pairs are constructed uniformly in
\S\ref{sec:families}, with Zsigmondy primes of the relevant torus
exponents as the obstruction primes. The finite exception manifest in
\cite{Repo} covers the Zsigmondy exceptions, small-field maximality exceptions,
the exceptional outer automorphisms of $A_6$, and the other isolated cases.
The manifest identifies the relevant simple groups.
The only ones above $1.05\times 10^{7}$ are the four
alternating groups $A_{11},A_{12},A_{13},A_{14}$, handled by machine
certificates, and the three Lie-type groups $\PSL(6,2)$, $\Omega^+(8,2)$,
and $\Sp(4,8)$, treated in the uniform proofs with substitute primes. In the interval
$5\times10^5\leq |S|\leq1.05\times10^{7}$, the $38$ groups
$\PSL(2,q)$ are handled by the uniform proof of Theorem~\ref{thm:psl2};
the remaining $13$ groups, the $47$ groups below $5\times10^5$, and every
sporadic group are handled by machine certificates (\S\ref{sec:machine}).

Three structural facts keep the family analysis manageable. First, when
the two exhibited classes consist of \emph{maximal} subgroups of $S$, all
hypotheses of the machinery are automatic and only stability and the
valuation inequality need checking. Second, for each twisted family used
below, the automorphism structure induces no nontrivial permutation of the
types of the relevant \emph{relative} (twisted) diagram; field and diagonal
automorphisms preserve those types. Thus the selected parabolic classes are
stable under every automorphism. The twisted families,
including the triality groups ${}^3D_4(q)$ and the Ree groups, are
consequently among the easiest cases.
Third, when a graph automorphism does fuse the relevant parabolic classes,
the corresponding \emph{flag} parabolics (intersections of fused maximal
parabolics) become poset-maximal among stable classes, and a structural
lemma (Lemma~\ref{lem:P}) shows they satisfy all needed hypotheses; this
is the wreath-product analogue of the ``novelty'' phenomenon for maximal
subgroups of almost simple groups.

\subsection*{Computer verification}
The machine-certified ingredients were produced with GAP~4.16.0
\cite{GAP}, using its simple-group iterators, the character table
library \cite{CTblLib} and \textsc{Atlas} data \cite{ATLAS}, in the
pinned computational environment described in \S\ref{sec:machine}.
The verification scripts used in the proof are deterministic and fail closed:
any skipped
case, soft failure, ambiguous identification, missing datum, version
mismatch, or byte-level certificate difference is fatal to
verification, and superseded exploratory scripts are not part of the
verification suite. Every finite certificate records class
positions and guard fingerprints as well as orders, the obstruction
prime, and valuations. The scripts and logs are archived in a
versioned release \cite{Repo}; \S\ref{sec:machine} specifies what is
claimed from the machine and how it is certified, and also records
the verification records behind the arithmetic of the family proofs of
\S\ref{sec:families}.

The structural arguments of \S\S\ref{sec:reduction}--\ref{sec:criterion}
and the universal arithmetic deductions used in \S\ref{sec:families} are
formalized in Lean~4 and Rocq and checked by their kernels. Lean checks the
quotient inheritance and the minimal-order core
of the reduction: uniqueness and characteristic
simplicity of the minimal normal subgroup, non-solubility, soluble quotient,
trivial centralizer, and faithfulness of the conjugation map
$G\to\Aut(N)$. It also checks the product-supplement intersection, exact
order, maximality, and non-conjugacy statements of
\S\ref{sec:machinery}, under the normalized coordinate-action hypotheses of
Conventions~\ref{conv:X}--\ref{conv:coord}. A further development checks the
subgroup-product/lower-divisibility bridge and the universal-in-$k$
arithmetic contradiction of \S\ref{sec:criterion}. Finally, Lean checks the
universal arithmetic deductions of the exact $34$-branch arithmetic manifest
behind the family proofs, conditional on the published group, Levi,
outer-automorphism, maximality, and Zsigmondy inputs identified by the source
rows.

The direct-power step has a separate Rocq/MathComp kernel check: every finite
characteristically simple non-soluble group is an internal direct product of
$k>0$ automorphic copies of one non-soluble, non-abelian simple subgroup $S$,
with $|N|=|S|^k$. The strengthened Rocq theorem also constructs an isomorphism
from $N$ to the explicit external coordinate product of those same factors.

Lean then kernel-checks reindexing any such nonempty finite coordinate type to
$\{1,\ldots,k\}$, producing exactly the group equivalence $N\cong S^k$ used
in the subsequent construction. From that equivalence the Lean development
constructs the coordinate factors and proves that every automorphism permutes
them; constructs the injective homomorphism
$\Aut(S^k)\to\Aut(S)\wr S_k$; proves transitivity of the ambient $G$-action
from minimal normality by an orbit-support argument; identifies the inverse
image of $\Inn(S)^k$ as $N$; and checks the coordinate normalization and the
divisibility $|G/N|\mid |X/\Inn(S)|^k k!$.

The three theorem signatures, called \texttt{producer}, \texttt{reindex}, and
\texttt{consumer} in the repository, are bound by a fail-closed formal
interface manifest and checker. No unproved group-theoretic lemma is inserted
at the Rocq--Lean interface. No single kernel checks the audited translation
from MathComp's subgroup isomorphism and external function product to
Lean's group equivalence and function product; that translation therefore
remains an explicit trusted interface. The two kernels and their pinned
libraries are trusted components as well. The formal coverage and interface
manifests, together with the stated trust assumptions, are part of
\cite{Repo}.

Generative AI was used in the research and preparation of this paper;
a full declaration appears before the bibliography.

\subsection*{Proof and verification roadmap}
The proof has five separately auditable layers; each row of
Table~\ref{tab:roadmap} states where the claim is proved and how it is
verified.

\begin{table}[htbp]
\small
\begin{tabular}{@{}p{0.34\textwidth}p{0.24\textwidth}p{0.32\textwidth}@{}}
\toprule
Claim & Where & Verification \\
\midrule
Minimal counterexample has unique minimal normal subgroup $S^k$ and
embeds in $X\wr S_k$ &
Prop.~\ref{prop:min}, Convs.~\ref{conv:X}--\ref{conv:coord} &
Elementary structural proof; kernel-checked in Lean and Rocq/MathComp
\\[2pt]
Normalizers $B_V=N_G(V^k)$ of stable classes are maximal,
non-conjugate, of order $t|V|^k$ &
Lemmas~\ref{lem:A}--\ref{lem:C}, \ref{lem:P} &
Structural proof; kernel-checked in Lean; parabolic input from
\cite{Carter} \\[2pt]
Divisibility criterion: one stable pair with
$\eta>\vp{p}(x)$ excludes $S^k$ for every $k\geq2$ &
Thm.~\ref{thm:D} &
Structural proof; kernel-checked in Lean \\[2pt]
Stable pairs exist for every infinite family &
\S\ref{sec:families}, Thms.~\ref{thm:psl2}--\ref{thm:graph} &
Structural proofs from published order formulas; the universal
arithmetic is kernel-checked in Lean and independently reproduced by two
implementations (\S\ref{sec:machine}) \\[2pt]
Sporadic groups, Tits group, and the finite base are excluded &
Props.~\ref{prop:base}, \ref{prop:sporadic} &
GAP certificates with coverage cross-checks; $38$ upper-range
$\PSL(2,q)$ entries are covered by Thm.~\ref{thm:psl2} \\
\bottomrule
\end{tabular}
\caption{Proof and verification roadmap.}\label{tab:roadmap}
\end{table}

Every machine-backed claim in \S\ref{sec:machine} is mapped in the
repository README to its script and committed log file.

\section{Reduction to socle analysis}\label{sec:reduction}

\begin{lemma}[Conjugation lemma]\label{lem:conj}
If $G=AB$ for subgroups $A,B$, then $G=A^xB^y$ for all $x,y\in G$.
\end{lemma}

\begin{proof}
Since $G=AB$, write $xy^{-1}=ab$ with $a\in A$, $b\in B$. Then, as
subsets of $G$,
\[
A^xB^y=x^{-1}Ax\,y^{-1}By=x^{-1}A(ab)By
 =x^{-1}(AB)\,y=x^{-1}Gy=G. \qedhere
\]
\end{proof}

Consequently, property $\mathrm{P}$ depends only on the conjugacy classes
of maximal subgroups: it says every pair of \emph{distinct classes} has
one (hence every) representative pair factorizing $G$.

\begin{lemma}\label{lem:quot}
Property $\mathrm{P}$ is inherited by quotients.
\end{lemma}

\begin{proof}
Maximal subgroups of $G/K$ are the images of the maximal subgroups of $G$
containing $K$; non-conjugate ones lift to non-conjugate ones, and
factorizations project.
\end{proof}

\begin{proposition}[Minimal counterexample structure]\label{prop:min}
Let $G$ have least order among the non-soluble groups with property
$\mathrm{P}$.
Then $G$ has a unique minimal normal subgroup $N=S^k$ with $S$ non-abelian
simple and $k\geq 1$, $C_G(N)=1$, $G/N$ is soluble, and $G$ embeds in
$\Aut(S)\wr S_k$ with $N$ mapping onto $\Inn(S)^k$ and $G$ acting
transitively on the $k$ coordinates.
\end{proposition}

\begin{proof}
Every proper quotient of $G$ has property $\mathrm{P}$
(Lemma~\ref{lem:quot}), hence is soluble by minimality of $|G|$. The
soluble radical $R$ of $G$ is trivial: otherwise $G/R$ is non-soluble
(as $R$ is soluble) with property $\mathrm{P}$, contradicting minimality.
So the socle of $G$ is a direct product of non-abelian minimal normal
subgroups. If $N_1\neq N_2$ were two of them, then $G/N_1$ is soluble
while $N_2\cong N_2N_1/N_1\leq G/N_1$ is non-soluble, a contradiction.
Hence there is a unique minimal normal subgroup $N=S^k$, $S$ non-abelian
simple. Now $C_G(N)\cap N=Z(N)=1$ since $N$ is a direct product of
non-abelian simple groups; if $C_G(N)$ were nontrivial it would contain
a minimal normal subgroup of $G$, necessarily $N$ by uniqueness,
contradicting $C_G(N)\cap N=1$. So $C_G(N)=1$, and the conjugation
action embeds $G$ into $\Aut(N)=\Aut(S)\wr S_k$ in the standard way;
minimal normality of $N$ forces transitivity on coordinates. Published
references for the two standard
structural inputs used here are \cite[Lemma~2.10, pp.~173--174]{ZhangShi2009}
for the automorphism group of a direct power and
\cite[p.~432]{LucchiniMorini2002} for the resulting transitive wreath
realization in the unique-minimal-normal-subgroup setting. Finally, $G/N$ is
soluble because it is a proper quotient.
\end{proof}

\begin{convention}[Coordinate closure]\label{conv:X}
For $G$ as in Proposition~\ref{prop:min} let $G_1$ be the stabilizer of
the first coordinate and $\pi_1\colon G_1\to\Aut(S)$ the projection. Set
$X:=\pi_1(G_1)\cdot\Inn(S)$, the \emph{coordinate closure}; by
coordinate-transitivity all coordinates give the same $X$ up to conjugacy,
and after conjugating inside $\Aut(S)\wr S_k$ we may and do assume
$G\leq X\wr S_k$ with $\Inn(S)\leq X\leq\Aut(S)$.

The normalization is achieved as follows, computing in the exponent
convention of Lemma~\ref{lem:conj} (actions composing left to right).
Fix $g_j\in G$ carrying coordinate $1$ to coordinate $j$,
with $g_1=1$, and let $a_j\in\Aut(S)$ be its $1\!\to\!j$ component. If
$g\in G$ carries coordinate $i$ to coordinate $j$ through the component
$c\in\Aut(S)$, then $g_i\,g\,g_j^{-1}$ stabilizes coordinate $1$ with
component $a_i\,c\,a_j^{-1}$, which therefore lies in
$\pi_1(G_1)\leq X$. So replacing $G$ by $\delta G\delta^{-1}$, where
$\delta:=(a_1,\dots,a_k)\in\Aut(S)^k$ (equivalently, conjugating by
$(a_1^{-1},\dots,a_k^{-1})$ in the exponent convention of
Lemma~\ref{lem:conj}), places every component of every element
of $G$ in $X$, and $a_1=1$ leaves $X$ itself unchanged. Since property
$\mathrm{P}$ and all criteria below are invariant under replacing $G$
by a conjugate under $(a,\dots,a)\in\Aut(S)^k$, statements may also be
proved after normalizing $X$ up to $\Aut(S)$-conjugacy.
\end{convention}

\begin{convention}[Coordinate action and arithmetic
notation]\label{conv:coord}
For \emph{coordinate} calculations we fix the following convention.
Write $c_g(z):=gzg^{-1}$ (left conjugation), let $S_i$ be the $i$-th
coordinate factor of $N\cong S^k$, and let $\sigma_g\in S_k$ be the
permutation determined by $c_g(S_j)=S_{\sigma_g(j)}$. There are
component automorphisms $a_{g,i}\in\Aut(S)$ such that
\[
\bigl(c_g(z)\bigr)_i=a_{g,i}\bigl(z_{\sigma_g^{-1}(i)}\bigr)
\qquad(z\in N,\ 1\leq i\leq k):
\]
the value arriving at coordinate $i$ comes from the
\emph{inverse-permuted} source coordinate $\sigma_g^{-1}(i)$. From
$c_{gh}=c_g\circ c_h$ these data compose by
\[
\sigma_{gh}=\sigma_g\sigma_h,\qquad
a_{gh,i}=a_{g,i}\,a_{h,\sigma_g^{-1}(i)} .
\]
After the normalization of Convention~\ref{conv:X}, every $a_{g,i}$
lies in $X$. Throughout, $c_g$ denotes left conjugation on elements, whereas
$A^g=g^{-1}Ag$ denotes conjugation of subgroups; thus
$A^g=c_{g^{-1}}(A)$. Write
\[
x:=|X/\Inn(S)|,\qquad t:=|G/N| .
\]
Then $t$ divides $x^k\,k!$.
\end{convention}

The following observation is not needed for the proof of
Theorem~\ref{thm:main} but explains why the pairs used later are of
``supplement'' type.

\begin{remark}[Top--supplement pairs always factorize]
If $M\supseteq N$ is maximal and $B$ is maximal with $B\not\supseteq N$,
then $BN=G$ ($BN$ is a subgroup strictly containing the maximal $B$),
and since $N\leq M$ we get $MB\supseteq NB=G$. Hence $G$ has property
$\mathrm{P}$ if and only if $G/N$ has property $\mathrm{P}$ and the product
of every pair of non-conjugate maximal subgroups not containing $N$ is $G$.
\end{remark}

\section{Product supplements}\label{sec:machinery}

Fix $G$ as in Conventions~\ref{conv:X} and~\ref{conv:coord}:
$N=S^k\leq G\leq X\wr S_k$,
$t=|G/N|$, $x=|X/\Inn(S)|$. For $V\leq S$ write $[V]$ for its
$S$-conjugacy class. An $S$-class is \emph{$X$-stable} if it is invariant
under the action of $X$,
i.e.\ $V^a$ is $S$-conjugate to $V$ for every $a\in X$ ($\Inn$-stability
being automatic).

\begin{definition}\label{def:poset}
$\Pos_X(S):=\{\,[V] : 1<V<S,\ N_S(V)=V,\ [V]\ X\text{-stable}\,\}$,
partially ordered by containment up to $S$-conjugacy
($[V]\leq[U]$ iff $V\leq U^s$ for some $s\in S$). Call $V$
\emph{normally saturating} if $\langle V^W\rangle=W$ for every overgroup
$V\leq W\leq S$.
\end{definition}

\begin{remark}\label{rem:maxauto}
If $V$ is a \emph{maximal} subgroup of $S$ and $[V]$ is $X$-stable, then
$[V]\in\Pos_X(S)$ (maximal subgroups of a non-abelian simple group are
self-normalizing), $[V]$ is a maximal element of the poset, and $V$ is
normally saturating (its only overgroups are $V$ and $S$, and
$\langle V^S\rangle=S$ by simplicity). Thus the lemmas below apply to
maximal classes with no further hypotheses.
\end{remark}

\begin{lemma}[Existence]\label{lem:A}
Let $[V]\in\Pos_X(S)$ and set $B_V:=N_G(V^k)$. Then
\begin{enumerate}
\item $B_V\cap N=N_S(V)^k=V^k$;
\item $B_VN=G$;
\item $|B_V|=t\cdot|V|^k$.
\end{enumerate}
\end{lemma}

\begin{proof}
(1) An element of $N=S^k$ normalizes $V^k$ iff each coordinate normalizes
$V$. (2) Frattini-type argument: let $g\in G$. In the coordinate-action
notation of Convention~\ref{conv:coord}, $(V^k)^g=c_{g^{-1}}(V^k)$, and
since every coordinate of $V^k$ is $V$,
\[
(V^k)^g=\prod_{i=1}^{k} a_{g^{-1},i}(V),
\]
each factor is the image of $V$ under a component automorphism
$a_{g^{-1},i}\in X$. By $X$-stability there are $s_i\in S$ with
$a_{g^{-1},i}(V)=V^{s_i}$. So $(V^k)^g$ is $N$-conjugate to
$V^k$: there is $n\in N$ with $(V^k)^{gn}=V^k$, i.e.\ $gn\in B_V$, so
$g\in B_VN$. (3) $|B_V|=|B_VN|\cdot|B_V\cap N|/|N|=t\,|V|^k$.
\end{proof}

\begin{lemma}[Maximality]\label{lem:B}
If in addition $[V]$ is a maximal element of $\Pos_X(S)$ and $V$ is
normally saturating, then $B_V$ is maximal in $G$.
\end{lemma}

\begin{proof}
Let $B_V\leq H\leq G$. Since $H\supseteq B_V$ and $B_VN=G$ we have
$HN=G$, so $H$ surjects onto $G/N$ and in particular acts
coordinate-transitively.

\emph{Step 1 (the coordinate components of $H$ recover $X$).} Let $H_1,G_1$ be the
coordinate-$1$ stabilizers. For $g\in G_1$ write $g=hn$ ($h\in H$,
$n\in N$); $n$ fixes every coordinate, hence so does $h$, so $h\in H_1$
and $\pi_1(G_1)=\pi_1(H_1)\cdot\pi_1(N\cap G_1)=\pi_1(H_1)\cdot\Inn(S)$.
Therefore $\pi_1(H_1)$ and $\Inn(S)$ together generate $X$, and a class
of subgroups of $S$ is $X$-stable iff it is stable under $\pi_1(H_1)$.

\emph{Step 2 (Goursat's lemma and saturation: $H\cap N$ is a product).} Let
$T:=H\cap N\supseteq V^k$, let $W_i:=\pi_i(T)\leq S$ be the coordinate
projections and $K_i:=T\cap S_i$ the coordinate kernels ($S_i$ the $i$-th
factor). Then $V\leq K_i$ (as $V^k\cap S_i=V$ in coordinate $i$),
$K_i\trianglelefteq W_i$, and $V\leq W_i$. Since $K_i$ is a normal
subgroup of $W_i$ containing $V$, it contains
$\langle V^{W_i}\rangle=W_i$ by saturation. So $K_i=W_i$ for every $i$
and $T=\prod_i W_i$.

\emph{Step 3 (the $W_i$ force $H=B_V$ or $H=G$).} $T$ is $H$-invariant.
For $h\in H_1$, projecting $T^h=T$ to the first coordinate gives
$W_1^{\pi_1(h)}=W_1$, so $[W_1]$ is $\pi_1(H_1)$-stable, hence
$X$-stable by Step~1. For any $i$, coordinate-transitivity of $H$
supplies $h\in H$ carrying coordinate $i$ to coordinate $1$ through a
component $c\in X$, and projecting $T^h=T$ gives $W_1=W_i^{\,c}$;
applying $X$-stability of $[W_1]$ to $c^{-1}$ gives $[W_i]=[W_1]$.
Write $[W]$ for this common $X$-stable class.
If $W=S$ then $T=N$, so $H\supseteq N$ and $H=HN=G$.
Otherwise $1<V\leq W<S$. Form the normalizer tower
$W\leq N_S(W)\leq N_S(N_S(W))\leq\cdots$; it is strictly increasing
until it becomes self-normalizing, and it cannot reach $S$ (the
penultimate term would be a proper normal subgroup of the simple group
$S$). Its terminal member $U$ is proper, self-normalizing, and $[U]$ is
$X$-stable ($N_S(W^a)=N_S(W)^a$, so stability propagates up the tower),
i.e.\ $[U]\in\Pos_X(S)$ with $[U]\geq[W]\geq[V]$. Maximality of $[V]$ in
$\Pos_X(S)$ forces $[U]=[V]$; then $|U|=|V|$ and $V\leq W\leq U$ gives
$V=W=U$. Since $V\leq W_i$ by Step~2 and $[W_i]=[W]$ gives
$|W_i|=|W|=|V|$, each $W_i=V$; so $T=\prod_i W_i=V^k$, hence
$|H|=|HN|\cdot|H\cap N|/|N|=t\,|V|^k=|B_V|$
and $H=B_V$.
\end{proof}

\begin{lemma}[Non-conjugacy]\label{lem:C}
Distinct $[V]\neq[W]$ in $\Pos_X(S)$ give non-conjugate $B_V,B_W$.
\end{lemma}

\begin{proof}
If $B_V^g=B_W$ then $(B_V\cap N)^g=B_W\cap N$ (as $N\trianglelefteq G$),
i.e.\ $(V^k)^g=W^k$. Reading off any coordinate, $W$ is the image of $V$
under an element of $X$ composed with a coordinate permutation, so
$[W]=[V^a]=[V]$ for some $a\in X$ by $X$-stability, a contradiction.
\end{proof}

The next lemma supplies, structurally, all hypotheses of
Lemma~\ref{lem:B} for the non-maximal (``novelty'') classes used in
\S\ref{sec:families}: the standard parabolic subgroups whose maximal
parabolic overgroups are fused by a graph automorphism. We use the
notation of \cite{Carter}. Here $S$ is an untwisted simple group of Lie
type with a split $BN$-pair of rank $\ell\geq2$; all our applications
are untwisted. Write $B=U\rtimes T$ for a Borel subgroup, with unipotent
radical $U$ and maximal torus $T$, and let $\Delta$ be the node set of
the Dynkin diagram. For $J\subseteq\Delta$, let $Q_J\supseteq B$ be the
standard parabolic whose Levi subgroup $L_J$ is generated by $T$ and the
root subgroups corresponding to $J$. The maximal parabolic omitting node
$i$ is $P_i=Q_{\Delta\setminus\{i\}}$.

\begin{lemma}[Parabolic novelties]\label{lem:P}
Fix $J\subseteq\Delta$ and put $V:=Q_J$. Suppose $X$ is such that
\begin{enumerate}
\item[(a)] the $S$-class $[Q_J]$ is $X$-stable, and
\item[(b)] for every $J\subsetneq K\subsetneq\Delta$ the class $[Q_K]$
is \emph{not} $X$-stable.
\end{enumerate}
Then $[V]\in\Pos_X(S)$, $[V]$ is a maximal element of $\Pos_X(S)$, and
$V$ is normally saturating. Consequently Lemmas~\ref{lem:A}--\ref{lem:C}
apply to $V$ with no computational hypothesis.
\end{lemma}

\begin{proof}
Parabolic subgroups are self-normalizing
\cite[Thm.~8.3.3, p.~112]{Carter}, so
$[V]\in\Pos_X(S)$ by (a).

\emph{Overgroups.} Every subgroup containing $V$ contains $B$, and every
subgroup containing a Borel subgroup is a parabolic $Q_K$, $K\supseteq J$,
of the same $BN$-pair \cite[Thm.~8.3.2, p.~112]{Carter}; moreover $Q_K$
and $Q_{K'}$ are $S$-conjugate iff $K=K'$
\cite[Thm.~8.3.3, p.~112]{Carter}.

\emph{Poset-maximality.} If $[U]\in\Pos_X(S)$ with $V\leq U^s<S$ for some
$s$, then $U^s=Q_K$ for some $J\subseteq K\subsetneq\Delta$; $K=J$ gives
$[U]=[V]$, while $J\subsetneq K$ contradicts (b), since $[U]=[Q_K]$ would
be $X$-stable.

\emph{Saturation.} Let $V\leq W\leq S$, so $W=Q_K$ ($K\supseteq J$) or
$W=S$; we must show $\langle V^W\rangle=W$. The normal closure of $V$ in
$W$ contains that of $B$. Write $W=U_K\rtimes L_K$ (Levi decomposition;
for $W=S$ read $U_K=1$, $L_K=S$). Then $U_K\leq B$, and modulo $U_K$ the
image of $B$ is a Borel subgroup $B_L=T\cdot(U\cap L_K)$ of $L_K$
containing the full torus $T$. For each node $\alpha\in K$ the
reflection representative $n_\alpha$ lies in $L_K$, and
$U_{-\alpha}=U_\alpha^{\,n_\alpha}$. The subgroup
$\langle B_L^{\,L_K}\rangle$ contains every $L_K$-conjugate of every
subgroup of $B_L$; in particular it contains $T$, each $U_\alpha$ for
$\alpha\in K$, and each $U_{-\alpha}=U_\alpha^{\,n_\alpha}$.
These generate $L_K$ by the root-subgroup description of the Levi
decomposition \cite[\S8.5, pp.~118--120]{Carter}. So
$\langle B^W\rangle\supseteq U_K\cdot L_K=W$.
\end{proof}

\section{The divisibility criterion}\label{sec:criterion}

For a prime $p$, $\vp{p}$ denotes the $p$-adic valuation and $s_p(k)$ the
digit sum of $k$ in base $p$, so that $\vp{p}(k!)=(k-s_p(k))/(p-1)$.

\begin{theorem}[Divisibility criterion]\label{thm:D}
Let $S,X$ be as in Conventions~\ref{conv:X} and~\ref{conv:coord}, and
let $[V]\neq[W]$ be
maximal elements of $\Pos_X(S)$, both normally saturating. If some prime
$p$ satisfies
\[
\eta=\eta_p(V,W):=\vp{p}(|S|)-\vp{p}(|V|)-\vp{p}(|W|)\;>\;\vp{p}(x),
\]
then for every $k\geq2$, no group $G$ with unique minimal normal
subgroup $S^k$ and coordinate closure $X$ has property $\mathrm{P}$.
\end{theorem}

\begin{proof}
By Lemmas~\ref{lem:A}--\ref{lem:C}, $B_V$ and $B_W$ are non-conjugate
maximal subgroups of $G$, so property $\mathrm{P}$ forces $G=B_VB_W$ and
\[
|B_V\cap B_W|=\frac{|B_V|\,|B_W|}{|G|}
 = t\cdot\Bigl(\frac{|V|\,|W|}{|S|}\Bigr)^{k}\in\Z ,
\]
whence $p^{\eta k}\mid t$. But
\[
\vp{p}(t)\;\leq\;k\,\vp{p}(x)+\vp{p}(k!)
 \;=\;k\,\vp{p}(x)+\frac{k-s_p(k)}{p-1}
 \;<\;k\bigl(\vp{p}(x)+1\bigr)\;\leq\;\eta k ,
\]
a contradiction.
\end{proof}

\begin{remark}\label{rem:crit}
(i) The criterion is uniform in $k$, and this uniformity handles the
infinite families. No fixed counting ratio can do so, because
$(k!)^{1/k}\to\infty$ and therefore eventually exceeds every fixed ratio
$|V||W|/|S|$. (ii) Only the \emph{existence} of
$B_V,B_W$ is used; no classification of the maximal subgroups of $G$
(or even of $S$) is needed. In particular the criterion is insensitive
to completeness of tabulated maximal-subgroup data: it needs only that
the two exhibited subgroups are maximal, or are novelties covered by
Lemma~\ref{lem:P}, and that their classes are stable. (Machine
\emph{verification} of stability for
the finite base does use enumerated subgroup data; see the trust
discussion in \S\ref{sec:machine}.)
\end{remark}

The criterion has a $k=1$ counterpart, which recovers the almost
simple case.

\begin{theorem}[$k=1$: the almost simple case]\label{thm:Dprime}
Let $S\leq G\leq\Aut(S)$ with $S$ non-abelian simple, and put $X:=G$,
so that $x:=|G/S|=t$. Let $[V]\neq[W]$ be maximal elements of
$\Pos_X(S)$, both normally saturating, and set $B_V:=N_G(V)$,
$B_W:=N_G(W)$. Then:
\begin{enumerate}
\item $B_V\cap S=V$, $B_VS=G$, and $|B_V|=t|V|$;
\item $B_V$ and $B_W$ are non-conjugate maximal subgroups of $G$;
\item if $\eta=\vp{p}(|S|)-\vp{p}(|V|)-\vp{p}(|W|)>\vp{p}(x)$ for some
prime $p$, then $G$ does not have property $\mathrm{P}$.
\end{enumerate}
\end{theorem}

\begin{proof}
\emph{Supplement and order.} By the Frattini argument via $X$-stability,
as in Lemma~\ref{lem:A}, we have $B_VS=G$. Moreover,
$B_V\cap S=N_S(V)=V$, and hence $|B_V|=t|V|$.

\emph{Maximality.} The proof of Lemma~\ref{lem:B} applies without the
Goursat step; normal saturation enters only through that step and is
not needed at $k=1$ (the hypothesis is kept for uniformity with
Theorem~\ref{thm:D}). In detail, let $B_V\leq H\leq G$ and
$T:=H\cap S$, an $H$-invariant subgroup containing $V$. If $T=S$ then
$H\supseteq S$ and $H=HS=G$. Otherwise, since $B_V\leq H$ and
$G=B_VS$, we have $G=HS$. As $T$ is
$H$-invariant, for $g=hs$ with $h\in H$ and $s\in S$ we have
$T^g=(T^h)^s=T^s$; hence $[T]$ is $G$-stable. Since $X=G$, it is
therefore $X$-stable, as in Step~1 of Lemma~\ref{lem:B}. The normalizer
tower of $T$ ends at a proper
self-normalizing subgroup $U$. Stability propagates up the tower, so
$[U]\in\Pos_X(S)$ and $[U]\geq[V]$. Poset-maximality of $[V]$ forces
$[U]=[V]$, and $V\leq T\leq U$ then gives $T=V$. Thus
$H\leq N_G(V)=B_V$.

\emph{Non-conjugacy.} If $B_V^g=B_W$ with $g\in G$, then
$V^g=(B_V\cap S)^g=B_W\cap S=W$, and $X$-stability of $[V]$ (here
$X=G$) gives $[W]=[V]$, a contradiction.

\emph{The arithmetic.} Property $\mathrm{P}$ forces $G=B_VB_W$, so
$t\,|V||W|/|S|\in\Z$, i.e.\ $p^\eta\mid t=x$, contradicting
$\eta>\vp{p}(x)$.
\end{proof}

\begin{convention}\label{conv:excluded}
From now on, ``$S$ is \emph{excluded}'' means: for every $X$ with
$\Inn(S)\leq X\leq\Aut(S)$ there exist classes as in
Theorem~\ref{thm:D}; consequently no finite group $G$ with unique
minimal normal subgroup $S^k$ ($k\geq2$, any admissible embedding) has
property $\mathrm{P}$. The almost simple branch $k=1$ is supplied separately
by Tikhonenko--Tyutyanov \cite{TT2010}.
\end{convention}

\section{The machine base}\label{sec:machine}

All machine certificates were produced in a pinned computational
environment: GAP~4.16.0 \cite{GAP} with
CTblLib~1.3.11 and AtlasRep~2.1.11, overlaid with the
\texttt{ContainedConjugates} fix from GAP commit
\texttt{b12f8342d641075d58fcbe62cc00dd433d7b8e18}; unpatched
GAP~4.16.0 is rejected. Both GAP global pseudorandom sources are
reset at process start to the recorded seed $1034$, so internal
library choices are reproducible; proof scripts contain no explicit random witness search.
The repository README in \cite{Repo} maps the short sweep labels used below
to the corresponding script and certificate filenames and maps every paper
claim to its verification artifacts.

\begin{proposition}\label{prop:base}
Every non-abelian simple group $S$ with $|S|\leq 1.05\times10^{7}$ is
excluded (in the sense of Convention~\ref{conv:excluded}).
\end{proposition}

\begin{proof}
We combine the uniform $\PSL(2,q)$ argument of
Theorem~\ref{thm:psl2}, whose proof of these cases does not appeal to
the present proposition, with the finite computation in GAP organized
as the two searches below. The group inventory is cross-checked against
GAP's \texttt{SimpleGroupsIterator} over two disjoint ranges, so that no
group in range is skipped: the $47$ groups with $|S|<5\times10^5$
(log \texttt{verify\_coverage.log}) and the $51$ groups with
$5\times10^5\leq|S|\leq1.05\times10^{7}$ (log
\texttt{verify\_coverage\_big.log}). Of the latter $51$, the $38$
$\PSL(2,q)$ entries are handled by Theorem~\ref{thm:psl2} and its arithmetic
verification rather than by a redundant GAP maximal-subgroup enumeration;
sweeps J3 and J6 exhaust the other $13$. The coverage check also confirms
that the only two cases in this upper range not excluded by maximal-class
pairs, $\Sp(4,4)$ at
$x=4$ and $\PSL(5,2)$ at $x=2$, are excluded by the novelty certificates
of item~(2).

\begin{enumerate}
\item \emph{Maximal-class pairs (sweeps J, J2, J4, J3, J5, J6).}
For each simple group $S$ treated computationally and each $X$ with
$\Inn(S)\leq X\leq\Aut(S)$, the program lists the $X$-stable classes of
maximal subgroups of $S$ and searches for pairs $([V],[W])$ and primes
$p$ satisfying the inequality of Theorem~\ref{thm:D}.

\emph{Certificate identity and coverage.} To identify coordinate closures
independently of GAP's internal class numbering, each closure is assigned a
SHA-256 fingerprint derived from its action on the pinned maximal-subgroup
classes, together with intrinsic quotient and kernel data. The checker verifies
that every required closure appears exactly once and that each certificate is
attached to one such closure. Each certificate records the selected classes,
their orders, the obstruction prime, and the valuation gap.
For all but six simple groups in range, a maximal-class pair is found for every
$X$. For $S=A_5$, for example, the logged tuple
$(|A_4|,|S_3|,5,1)$ excludes $A_5^k$ for every $k$ and every $X$. In
addition, the alternating groups $A_{11}$--$A_{14}$, whose orders
exceed the base bound, carry certificates of the same form for every
$X$ (sweep J5, log \texttt{sweepJ5\_smallAn.log}); Theorem~\ref{thm:an}
uses these for $11\leq n\leq14$.
\item \emph{Novelty pairs (sweeps K, K2, K3, K4).} Six simple groups,
for the indicated coordinate closures, admit too few stable maximal
classes because an outer automorphism fuses classes: $\PSL(3,2)$ and $\PSL(2,11)$ at
$x=2$; $A_6$ for $X/\Inn(S)\in\{2_2,2_3,2^2\}$; $\PSL(3,4)$ for several
$X$; $\Sp(4,4)$ at $x=4$; and $\PSL(5,2)$ at $x=2$. Each is excluded by a
pair involving non-maximal classes that are maximal in $\Pos_X(S)$. For
$\PSL(3,2)$, $\PSL(2,11)$, $A_6$, and $\PSL(3,4)$, normal saturation is
verified computationally by sweep K2; for $\Sp(4,4)$ and $\PSL(5,2)$,
it follows structurally from Lemma~\ref{lem:P}.

\emph{Normal-saturation check and trust boundary.} For the four computational
novelty cases, three independent procedures enumerate the relevant overgroups
and embedding orbits and verify normal saturation; any disagreement causes
verification to fail. Each selected subgroup class is tied to its fingerprint
and normal-closure result, so same-order non-conjugate classes are not
identified by order alone. These calculations assume the completeness of GAP's
pinned subgroup-class enumerations. This assumption is cross-checked by the
independent procedures and, for the sporadic groups, against the published
\textsc{Atlas} data.
Sample certificates: $\PSL(3,2)$ at
$x=2$ is excluded by the pair $(S_3,\,7{:}3)$ at $r=2$ with logged
valuation gap $2>1$. The case $\Sp(4,4)$ at $x=4$ is excluded by the pair (Borel,
subfield $\Sp(4,2)$) at $r=5$ and $17$. The case $\PSL(5,2)$ at $x=2$
is excluded by the
incident-flag parabolics $(Q_{\{2,3\}},Q_{\{1,4\}})$ (the
point--hyperplane and line--$3$-space flag stabilizers, in the standard
parabolic notation of Lemma~\ref{lem:P}) at $r=5$ and $31$. The last
two are
instances of the uniform constructions of Theorem~\ref{thm:graph} below.
\end{enumerate}
\end{proof}

\begin{proposition}\label{prop:sporadic}
Every sporadic simple group is excluded, as is the Tits group
${}^2F_4(2)'$.
\end{proposition}

\begin{proof}
Six of the $26$ sporadic groups ($M_{11}$, $M_{12}$, $M_{22}$,
$M_{23}$, $J_1$, $J_2$) have order at most $1.05\times10^{7}$ and fall
under Proposition~\ref{prop:base}. The remaining $20$ ($M_{24}$, $J_3$,
$J_4$, $\mathrm{HS}$, $\mathrm{McL}$, $\mathrm{Co}_1$, $\mathrm{Co}_2$,
$\mathrm{Co}_3$, $\mathrm{Suz}$, $\mathrm{He}$, $\mathrm{Ru}$, $\mathrm{O'N}$,
$\mathrm{Fi}_{22}$, $\mathrm{Fi}_{23}$, $\mathrm{Fi}_{24}'$, $\mathrm{HN}$,
$\mathrm{Ly}$, $\mathrm{Th}$, $B$, $M$) and the Tits group are handled by a dedicated
computation (sweep~M; script \texttt{sweepM\_sporadic.g}, certificate
\texttt{sweepM\_sporadic.log}): stable maximal classes and pair searches use the
\textsc{Atlas} maximal-subgroup data via the GAP character table
library \cite{CTblLib,ATLAS} in the pinned environment above.
Every record gives the character-table identifier and
\texttt{Maxes} list position. If $\Out(S)=1$, stability is automatic. If
$|\Out(S)|=2$, the program accepts a class only when its order occurs exactly
once in the complete pinned \texttt{Maxes} list; automorphisms preserve order,
so that class is stable. No stored class fusion is claimed or used.
Each of the resulting $42$ selected class records is also bound, by
group, list position, identifier, and order, to a published maximal-subgroup
record. For the non-Monster sporadics these are Wilson's corrected tables
\cite[\S4, pp.~66--68]{Wilson2017}, whose published-proof provenance is
summarized in \cite[\S2, pp.~23--24]{BreuerMalleOBrien2017}; the Tits entries
are covered independently by \cite[pp.~553--563]{Wilson1984} and
\cite{Tchakerian1986}; and the Monster entries are given by
Theorems~1.1--1.2 and Table~1 of \cite[pp.~862--863]{DietrichEtAl2026}.
Sample certificates: the Monster is excluded via positions $45,46$ by
the pair $(59{:}29,\,41{:}40)$ at $r=71$, and the Baby Monster via
positions $29,30$ by the pair
$(L_2(11){.}2,\,47{:}23)$ at $r=31$. By Remark~\ref{rem:crit},
completeness is logically needed here only for the
unique-order stability argument when $|\Out(S)|=2$. The source map is
machine-checked against the certificate rather than maintained only in prose.
\end{proof}

Beyond the base, the family proofs of \S\ref{sec:families} carry their
own numeric verification records: the valuation inequalities of
Theorems~\ref{thm:psl2}, \ref{thm:an},
and~\ref{thm:twisted2}--\ref{thm:graph} were verified mechanically over
large parameter ranges ($1272$ prime powers for Theorem~\ref{thm:psl2};
$n\leq10^4$ for Theorem~\ref{thm:an}; $7892$ parameter instances,
ranks up to $25$ and fields up to $3000$, for
Theorems~\ref{thm:twisted2}--\ref{thm:graph}), with zero failures and
with the Zsigmondy-exception set as documented in the proofs;
the Levi-exponent inequalities behind Theorem~\ref{thm:nograph} are
checked separately by the independent symbolic verifier in \cite{Repo}.
These finite-range runs are regression checks, not proofs of the statements
quantified over all ranks and fields.

The release also contains an exact $34$-branch arithmetic manifest for
Theorems~\ref{thm:psl2}--\ref{thm:graph}. A universal checker matches every
branch to its cited order formula and Zsigmondy source, verifies the stated
exponent bounds for the full parameter ranges, and reproduces the treatment of
the exceptional cases. The positive-base Zsigmondy exception occurs in twelve
branches, corresponding to seven distinct exception records; each is checked
against the substitute-prime and finite certificates. A
separate symbolic implementation independently reproduces the $34$ branches,
bounds, residue arguments, and exceptional valuations without reading the
manifest. The same universal deductions are kernel-checked in Lean; the
published group, Levi, outer-automorphism, maximality, and Zsigmondy results
remain cited inputs rather than newly formalized classification theorems.

\section{The infinite families}\label{sec:families}

Throughout this section $S$ is simple, $X$ is any group with
$\Inn(S)\leq X\leq\Aut(S)$, and $x=|X/\Inn(S)|$. By
Remark~\ref{rem:maxauto}, when the two exhibited classes consist of
maximal subgroups of $S$, only their $X$-stability and the valuation
inequality require proof. We use Zsigmondy's theorem
\cite[specialized theorem, p.~283]{Zsigmondy}
throughout: for integers $a\geq2$, $e\geq3$ with $(a,e)\neq(2,6)$, there
is a prime $r$ with $\ord_r(a)=e$ (a \emph{primitive prime divisor} of
$a^e-1$); such $r$ satisfies $r\equiv1\pmod e$. For $q=p^f$ we write
$\Phi_e=\Phi_e(q)$ for cyclotomic values and always apply Zsigmondy to
the prime $p$ with exponent $ef$, so $\ord_r(p)=ef$ and
$r\equiv1\pmod{ef}$. The source's additional positive-base exception has
exponent $2$, which lies outside every invocation here and therefore does not
apply. An invocation-by-invocation audit of the hypotheses and exceptions
is part of \cite{Repo}. The $b=1$ specialization and both exceptions are
independently restated in \cite[Thm.~2.2, p.~2]{Jones2007}.

The simple-group and outer-automorphism order formulas used below are
those of \cite[Tables~5--6, p.~xvi]{ATLAS}; for the classical families,
the same formulas and their central divisors are also displayed in
\cite[\S\S2.1--2.4, pp.~x--xii]{ATLAS}. As an independent published
derivation and tabulation we use
\cite[\S10, pp.~220--221, and summary table, p.~239]{Carter1965}.
For parabolics, the statement that the unipotent radical is extended by
the group obtained from the complementary subdiagram (with node-orbits in
the twisted case) is \cite[\S3.6, p.~xv]{ATLAS}; the Levi decomposition
and its root-subgroup description are
\cite[\S8.5, pp.~118--120]{Carter}. Thus the cyclotomic bounds below are
not inferred from numeric tests: they come from the displayed group-order
formulas applied to the named Levi types. The formula-to-source and
parameter audit, including each exceptional substitute, is archived in
\cite{Repo}. In the exceptional case over $\F_2$ below,
$\Omega^+(8,2)$ and $P\Omega^+(8,2)$ denote the same simple group; we use
the shorter former notation when naming that finite group. For classical
groups, argument notation such as $\Sp(4,2^f)$ denotes a named simple group,
whereas subscript notation such as $\Sp_2(q)$ and $\GL_2(q)$ is used for Levi
factors.

\subsection{Projective special linear groups of rank one}

\begin{theorem}\label{thm:psl2}
$S=\PSL(2,q)$, $q=p^f\geq4$, is excluded.
\end{theorem}

\begin{proof}
For $q\in\{4,5,7,8,9,11\}$ this is machine-certified by
the lower-range J/K certificates of \S\ref{sec:machine}. All other simple parameters are covered
uniformly below: odd $q\geq13$ by Case~A and even $q=2^f$, $f\geq3$, by
Case~B. Recall
$|S|=q(q^2-1)/d$ with $d=\gcd(2,q-1)$, and
$\Out(S)=C_d\times C_f$, so $x\mid df$: that is, $x\mid2f$ for $q$ odd
and $x\mid f$ for $q$ even.

\emph{Case A: $q$ odd, $q\geq13$.} Let $U:=N_S(T_+)$ and $V:=N_S(T_-)$
be the normalizers of the split and nonsplit maximal tori. These are dihedral
groups of orders $q-1$ and $q+1$, respectively. Both subgroups are maximal in $S$
for $q\geq13$ by Dickson's theorem \cite{Dickson,Huppert}; see also
\cite[Tables~8.1--8.2, p.~376, and Table~8.7, p.~380]{BHRD}.
All split maximal tori are $S$-conjugate, as are all nonsplit maximal
tori, and every automorphism preserves each type (there is no
graph automorphism in type $A_1$; diagonal and field automorphisms act
algebraically), so $[U]$ and $[V]$ are $X$-stable for every $X$. Now
$|U||V|/|S|=d/q$, so for the defining prime $p$:
$\eta=\vp{p}(|S|)-0-0=f$, while $x\mid2f$ gives
$\vp{p}(x)\leq\vp{p}(f)\leq\log_pf<f=\eta$. Theorem~\ref{thm:D} applies.

\emph{Case B: $q=2^f$, $f\geq3$.} Here $|S|=q(q^2-1)$ and $x\mid f$.
Let $U$ be a Borel subgroup (the normalizer of a Sylow $2$-subgroup), of
order $q(q-1)$, and let $V:=N_S(T_+)$ be dihedral of order $2(q-1)$.
Both $U$ and $V$ are maximal for $q\geq8$ \cite{Dickson,Huppert}; see also
\cite[Tables~8.1--8.2, p.~376]{BHRD}. $[U]$ is the class of Sylow-$2$
normalizers, and $[V]$ is the unique class of split torus normalizers;
both are $X$-stable. Now $|U||V|/|S|=2(q-1)/(q+1)$ with $q+1$ odd and coprime to
$2(q-1)$, so it suffices to find a prime $r\mid q+1$ with
$\vp{r}(q+1)>\vp{r}(f)$. If $f\neq3$, take $r$ a primitive prime divisor
of $2^{2f}-1$: then $r\mid(2^{2f}-1)/(2^f-1)=q+1$ and
$r\equiv1\pmod{2f}$, so $r>f$ and
$\vp{r}(x)=0<\vp{r}(q+1)=\eta$. If $f=3$, then $q=8$ and
$q+1=9$; taking $r=3$ gives $\eta=2>1\geq\vp{3}(x)$.
\end{proof}

\subsection{Alternating groups}

\begin{theorem}\label{thm:an}
$S=A_n$, $n\geq5$, is excluded.
\end{theorem}

\begin{proof}
For $n\leq14$ this is machine-certified: the cases $n\leq10$ fall under
Proposition~\ref{prop:base} (including $n=6$, whose exceptional outer
automorphisms fuse classes; all five $X$ are handled there), and
the cases $11\leq n\leq14$ are covered by the sweep-J5 certificates of
\S\ref{sec:machine}.
Assume $n\geq15$, so $\Aut(A_n)=S_n$ and $x\mid2$.

Let $p$ be the largest prime $\leq n$. We first record the bound
$p\geq n/2+3$. By Nagura's theorem
\cite[Theorem, pp.~180--181]{Nagura}, there is a prime in $(y,6y/5]$
for $y\geq25$. Thus, for $n\geq31$, there is a prime in $(5n/6,n]$,
and $5n/6>n/2+3$ for $n>18$. For $15\leq n\leq30$, the largest prime
$p\leq n$ satisfies the bound by inspection:
$p=13,13,17,17,19,19,19,19,23,23,23,23,23,23,29,29$ for
$n=15,\dots,30$. In particular, $n/2<p\leq n$.

For $m<n/2$ let
$M_m:=(S_m\times S_{n-m})\cap A_n$, and for $n$ even let
$I:=(S_{n/2}\wr S_2)\cap A_n$. Each $M_m$ is maximal in $A_n$, and $I$
is maximal except in a finite list of cases, all with $n\leq14$
\cite[Thm.~1 and Tables~I--II, pp.~366--368]{LPS1987}.
Conjugation by $S_n$ preserves the partition shape
(resp.\ block structure), each shape is a single $A_n$-class, and
$\Aut(A_n)=S_n$ for $n\neq6$: all these classes are $X$-stable for every
$X$.

For any two distinct classes among
$\{M_m:\,n-p<m<n/2\}\cup\{I\ \text{if $n$ even and}\ n/2<p\}$, we have
$\vp{p}(|A_n|)=1$, while $\vp{p}$ of each
subgroup order is $0$ (all parts of the partitions involved are $<p$),
and $\vp{p}(x)\leq\vp{p}(2)=0$. Thus $\eta=1>0$, and
Theorem~\ref{thm:D} applies, provided two such classes
exist. The number of admissible intransitive $m$ is
$p-\lfloor n/2\rfloor-1$, plus one for $I$ when $n$ is even; two classes
exist because $p\geq n/2+3$.
\end{proof}

\subsection{Lie type without graph symmetry}

\begin{theorem}\label{thm:nograph}
Let $S$ be a simple group of Lie type of rank $\geq2$ over $\F_q$,
$q=p^f$, of a type with no symmetry of its Dynkin diagram in the given
characteristic:
\[
C_n=\PSp(2n,q)\ (n\geq2,\ p\ \text{odd if}\ n=2),\quad
B_n=\Omega(2n{+}1,q)\ (n\geq3,\ q\ \text{odd}),
\]
\[
G_2(q)\ (p\neq3),\quad F_4(q)\ (p\neq2),\quad E_7(q),\quad E_8(q).
\]
Then $S$ is excluded.
\end{theorem}

\begin{proof}
Write $h$ for the relevant torus exponent: $h=2n$ for $B_n/C_n$, and
$h=6,12,18,30$ for $G_2,F_4,E_7,E_8$.

Let $P_1,P_2$ be maximal parabolic subgroups from two distinct classes
(rank $\geq2$ guarantees these; take the ends of the diagram). They are
maximal subgroups of $S$: any overgroup of a parabolic contains a Borel
subgroup, and every subgroup containing a Borel is parabolic
\cite[Thm.~8.3.2, p.~112]{Carter}.

\emph{Stability.} $\Aut(S)$ is generated by inner, diagonal, field and
graph automorphisms \cite{Steinberg},
\cite[Def.~3.1 and Thm.~3.4, pp.~31--32]{BrotoMollerOliver2019}. Diagonal
and field automorphisms preserve types of vertices of the building,
hence each parabolic class; type-permuting automorphisms arise only from
diagram symmetries, absent for the listed types in the listed
characteristics (the exceptional symmetries of $B_2/C_2$, $F_4$ in
characteristic~$2$ and $G_2$ in characteristic~$3$ are excluded by
hypothesis and treated in Theorem~\ref{thm:graph}). So every parabolic
class is $X$-stable for every $X$.

\emph{The obstruction prime.} Let $r$ be a primitive prime divisor of
$p^{hf}-1$. Then:
(i) $r\mid|S|$, since the order formula of each listed type contains the
factor $q^h-1$ (equivalently $\Phi_h(q)$)
\cite[Table~6, p.~xvi]{ATLAS}.
(ii) $r$ divides no maximal parabolic order: $|P_i|=q^{N_i}|L_i|
(q-1)^{c_i}$ with $L_i$ the semisimple part of the Levi factor, whose
order is a product of a $q$-power and factors $q^j-1$ with $j\leq h_i$,
$h_i$ the largest torus exponent of $L_i$; in every listed case every
proper Levi has $h_i<h$ ($B_n/C_n$: Levi types
$A_{i-1}\times C_{n-i}$ give $j\leq\max(i,2(n-i))<2n$; $G_2$: Levis
$A_1$, $j\leq2$; $F_4$: Levis $B_3,C_3,A_2{\times}A_1$, $j\leq6$; $E_7$:
Levis $E_6,D_6,A_6,\dots$, exponents $\leq12$; $E_8$: Levis
$E_7,D_7,A_7$, exponents $\leq18$). By primitivity $r$ divides none of
these, nor $q$, nor $q-1$.
(iii) $r\nmid x$: $r\equiv1\pmod{hf}$ with $h\geq4$, so
$r>hf\geq4f\geq d_S f\geq x$, where $d_S\leq\gcd(2,q-1)^{\,2}\leq4$
bounds the diagonal part and there is no graph part.

Thus $\eta=\vp{r}(|S|)\geq1>0=\vp{r}(x)$ for the pair $([P_1],[P_2])$, and
Theorem~\ref{thm:D} applies, provided $r$ exists.

\emph{Zsigmondy exceptions.} $p^{hf}=2^6$ forces $p=2$ and $hf=6$:
either $G_2(2)$ (not simple; $G_2(2)'\cong\PSU(3,3)$ is in the machine
base) or $\Sp(6,2)$ (machine base). For $h\in\{12,18,30\}$, $hf=6$ is
impossible; $B_n$ has $q$ odd.
\end{proof}

\subsection{Twisted groups of twisted rank at least two}

\begin{theorem}\label{thm:twisted2}
Let $S$ be one of
\[
\begin{gathered}
\PSU(n,q)\ (n\geq4),\quad
P\Omega^-(2n,q)={}^2D_n(q)\ (n\geq4),\\
{}^3D_4(q),\quad {}^2E_6(q),\quad
{}^2F_4(q)\ (q=2^f,\ f\ \text{odd}\geq3).
\end{gathered}
\]
Then $S$ is excluded.
\end{theorem}

\begin{proof}
Each of these groups has a twisted $BN$-pair of rank at least $2$, and
therefore has at least two classes of maximal parabolic subgroups $P_1,P_2$.
These subgroups are maximal in $S$, by the same argument as
in Theorem~\ref{thm:nograph}: every finite group of Lie type in defining
characteristic has a split $BN$-pair
\cite[Example~1.16 and \S1.3.5, pp.~10--11]{Hiss2011}, and Carter's
overgroup and self-normalizer theorems are stated for an arbitrary group
with a $BN$-pair \cite[Thms.~8.3.2--8.3.3, p.~112]{Carter}.

\emph{Stability.} We claim every parabolic class is $X$-stable for every
$X$. Every $a\in\Aut(S)$
permutes the parabolic subgroups: they are exactly the overgroups of the
normalizers of the Sylow $p$-subgroups (Borel subgroups), an
automorphism-invariant description; hence $a$ permutes the finitely many
parabolic classes, preserving inclusion up to conjugacy, i.e.\ acts on
the type set of the twisted building. Now use the automorphism description of
\cite[Def.~3.3 and Thm.~3.4, p.~32]{BrotoMollerOliver2019}, which sets the
separate graph factor equal to $1$ for a twisted group. The inner-diagonal and
Frobenius representatives
centralizing the defining Steinberg endomorphism preserve each orbit of
simple roots, hence each node of the relative (twisted) diagram. Inner
automorphisms fix classes, while the standard outer representatives normalize
the standard Borel and torus and therefore preserve the class of each relative
type. Thus they map every standard parabolic $Q_J$ to a conjugate of the same
type. Hence every automorphism fixes every
parabolic class. (For the rank-two group ${}^2F_4(q)$ one can also see
directly that the two maximal parabolic classes cannot be interchanged:
their Levi factors ${}^2B_2(q)$ and $\SL_2(q)$ are non-isomorphic. Note
that a diagram-asymmetry argument alone would not suffice at twisted
rank two: the twisted diagrams of $\PSU(4,q)$ and $\PSU(5,q)$ have a
symmetric underlying graph, and the untwisted example $\Sp(4,2^f)$
shows that a multiple bond by itself does not prevent a type-swapping
automorphism.)

\emph{The obstruction prime.} For each type, choose the exponent $e$ and
let $r$ be a primitive prime divisor of $p^{e}-1$. In each case $r$ divides
$|S|$ but neither $|P_1|$ nor $|P_2|$:
\begin{itemize}
\item $\PSU(n,q)$, $n$ odd: $e=2nf$, so $r\mid\Phi_{2n}(q)\mid q^n+1$,
which divides $|\SU(n,q)|=q^{n(n-1)/2}\prod_{i=2}^n(q^i-(-1)^i)$. The
Levi factor of $P_i$ ($i=1,2$) is a central product of $\GL_i(q^2)$ and
$\GU(n-2i,q)$: its $p'$-order is a product of factors $q^{2j}-1$
($j\leq2$) and $q^j-(-1)^j$ ($j\leq n-2$), all dividing $p^m-1$ with
$m\leq\max(4f,2(n-2)f)<2nf$.
\item $\PSU(n,q)$, $n$ even: $e=2(n-1)f$,
$r\mid\Phi_{2(n-1)}(q)\mid q^{n-1}+1$ (the $i=n-1$ factor). Levi
exponents are $\leq\max(4f,2(n-3)f,(n-2)f)<2(n-1)f$ for $n\geq4$.
\item ${}^2D_n(q)$, $n\geq4$: $e=2nf$, $r\mid\Phi_{2n}(q)\mid q^n+1$,
which divides $q^{n(n-1)}(q^n+1)\prod_{i=1}^{n-1}(q^{2i}-1)$.
The Levi factor of $P_i$ has type
$\GL_i(q)\times O^-(2(n-i),q)$, and its relevant exponents are at most
$\max(2f,2(n-1)f)<2nf$.
\item ${}^3D_4(q)$: $e=12f$, $r\mid\Phi_{12}(q)=q^4-q^2+1$, which
divides $q^8+q^4+1=(q^4+q^2+1)(q^4-q^2+1)$. Since
$|{}^3D_4(q)|=q^{12}(q^8+q^4+1)(q^6-1)(q^2-1)$, we have
$r\mid|S|$. The two Levi factors
are $\SL_2(q^3)\cdot(q-1)$ and $\SL_2(q)\cdot(q^3-1)$ (mod centers):
$p'$-orders divide $(q^6-1)(q-1)$ and $(q^2-1)(q^3-1)$, exponents
$\leq6f<12f$.
\item ${}^2E_6(q)$: $e=18f$, $r\mid\Phi_{18}(q)\mid q^9+1$, which
divides $q^{36}(q^{12}-1)(q^9+1)(q^8-1)(q^6-1)(q^5+1)(q^2-1)$. The four
maximal parabolic classes correspond to the four $\tau$-orbits of nodes
of $E_6$ ($\tau$ the order-two symmetry): $\{2\}$, $\{4\}$, $\{1,6\}$,
$\{3,5\}$; the complementary $\tau$-stable subdiagrams are $A_5$,
$A_2{\times}A_1{\times}A_2$, $D_4$, $A_2{\times}A_1{\times}A_1$ with induced
twists, so the Levi derived groups are $\SU(6,q)$,
$\SL_2(q){\times}\SL_3(q^2)$, ${}^2D_4(q)$-type, and
$\SL_3(q){\times}\SL_2(q^2)$, with cyclotomic exponents
$\leq10f,\,6f,\,8f,\,4f$ respectively, all $<18f$. Any two classes
serve as $P_1,P_2$. This diagram derivation agrees with the published list
$({}^2A_5(q),{}^2D_4(q),A_1(q){\times}A_2(q^2),
A_1(q^2){\times}A_2(q))$ in
\cite[Lemma~6.8, p.~82]{Montinaro2024}.
\item ${}^2F_4(q)$, $f$ odd $\geq3$: $e=12f$,
$r\mid\Phi_{12}(q)\mid q^6+1$, which divides
$|{}^2F_4(q)|=q^{12}(q^6+1)(q^4-1)(q^3+1)(q-1)$. The two Levi derived
groups are ${}^2B_2(q)\cong\Sz(q)$ and $\SL_2(q)$
\cite[Thm.~4.5(i)--(ii), p.~166]{Wilson2009} and
\cite[Main Theorem, pp.~52--53]{Malle1991}: $p'$-orders divide
$(q^2+1)(q-1)^2$ and
$(q^2-1)(q-1)$, exponents $\leq4f<12f$. (${}^2F_4(2)$ is not simple;
its derived group, the Tits group, is in the machine base.)
\end{itemize}

\emph{$r\nmid x$.} In every case $x$ divides $d_S\cdot c f$, where the
diagonal order $d_S$ divides $q+1$ (types ${}^2A$, ${}^2E_6$) or $4$
(type ${}^2D$) or is $1$ (${}^3D_4$, ${}^2F_4$), and the field part has
order $cf$ with $c\leq3$. Since $r\equiv1\pmod e$ with $e\geq6f$, $r$
exceeds the field-part order. If $r\mid d_S$, then either $r\mid q+1$,
which would give $\ord_r(p)\mid2f<e$, or $r\leq4$, contradicting
$r>e\geq6$. Hence
$\eta=\vp{r}(|S|)\geq1>0=\vp{r}(x)$, and Theorem~\ref{thm:D} applies.

\emph{Zsigmondy exceptions.} $r$ fails to exist only if $p^e=2^6$. Here
$e\geq6$ always, and $e=6$ occurs only for $\PSU(4,q)$ with $f=1$; then
$p^6=2^6$ forces $S=\PSU(4,2)\cong\PSp(4,3)$
\cite[Thm.~5.3, p.~61]{CameronNotes}; this group, of order $25920$, is
settled by the lower-range certificate \texttt{Sp43}. No other case
reaches $p^e=2^6$.
\end{proof}

\subsection{Twisted groups of twisted rank one}

\begin{theorem}\label{thm:twisted1}
Let $S$ be one of
\[
\PSU(3,q)\ (q\geq3),\quad
\Sz(q)={}^2B_2(q)\ (q=2^f,\ f\ \text{odd}\geq3),\quad
{}^2G_2(q)\ (q=3^f,\ f\ \text{odd}\geq3).
\]
Then $S$ is excluded.
\end{theorem}

\begin{proof}
Each group has a single class of parabolic subgroups, namely the Borel
subgroups $B=N_S(Q)$, $Q\in\Syl_p(S)$; $B$ is a \emph{maximal}
subgroup (rank-one $BN$-pair) and its class is $X$-stable for every $X$.
We pair it with a second $X$-stable maximal class avoiding a Zsigmondy
prime.

\emph{$\PSU(3,q)$, $q\geq8$.} $|S|=q^3(q^3+1)(q^2-1)/d$,
$d=\gcd(3,q+1)$; $|B|=q^3(q^2-1)/d$. Let $V$ be the stabilizer of a
nonisotropic point of the unitary geometry. It has order
$q(q-1)(q+1)^2/d$ and is maximal for $q\geq3$, $q\neq5$
\cite[Table~8.5, p.~379]{BHRD} (class $\mathcal{C}_1$). These stabilizers
form a single class, preserved by every semilinear automorphism.
Take $r$ a primitive prime divisor of $p^{6f}-1$. It exists because
the exceptional equation $p^{6f}=2^6$ would force $q=2$, for which $S$
is not simple. Then $r\mid\Phi_6(q)=q^2-q+1$,
which divides $(q^3+1)/(q+1)$, and hence $r\mid|S|$;
$\ord_r(p)=6f$ rules out
$r$ dividing $q(q^2-1)$ or $q(q-1)(q+1)^2$. Also $x\mid2df$, and
$r\equiv1\pmod{6f}$ gives $r>2f$, while $r\mid d\mid q+1$ would force
$\ord_r(p)\mid2f$. Thus $\eta:=\vp{r}(|S|)\geq1>0=\vp{r}(x)$.
For $q\in\{3,4,5,7\}$ (in particular the maximality exception $q=5$),
$S$ is in the machine base.

\emph{$\Sz(q)$.} $|S|=q^2(q^2+1)(q-1)$; $|B|=q^2(q-1)$. Let
$V:=D_{2(q-1)}$ be the normalizer of a cyclic subgroup of order $q-1$.
It is maximal by Suzuki's classification \cite{Suzuki1962}; see also
\cite[Thm.~7.3.5, p.~367, and Table~8.16, p.~384]{BHRD}. These
normalizers form a single class, stable under $\Aut(S)=S\rtimes C_f$. Let
$r$ be a primitive prime divisor of $2^{4f}-1$. Such a prime exists
because $4f\geq12$. Then $r\mid\Phi_4(q)=q^2+1$, so $r\mid|S|$,
$r\nmid|B|$, $r\nmid|V|$; $x\mid f$, and
$r\equiv1\pmod{4f}$ gives $r>f$. Thus
$\vp{r}(|S|)\geq1>0=\vp{r}(x)$.

\emph{${}^2G_2(q)$.} $|S|=q^3(q^3+1)(q-1)$; $|B|=q^3(q-1)$. Let
$V:=C_S(\iota)\cong\langle\iota\rangle\times\PSL(2,q)$ for an involution
$\iota$. The involutions form a single $S$-class, so $[V]$ is
$\Aut(S)$-stable. Moreover, $|V|=q(q^2-1)$, and $V$ is maximal
\cite[Thm.~C, pp.~33--34]{Kleidman1988}. Let $r$ be a primitive prime
divisor of $3^{6f}-1$. Such a prime exists because $6f\geq18$. Then
$r\mid\Phi_6(q)=q^2-q+1$, which divides $(q^3+1)/(q+1)$, and hence
$r\mid|S|$;
$\ord_r(3)=6f$ rules out $r\mid q^3(q-1)$ and $r\mid q(q^2-1)$;
$x\mid f<r$. So $\vp{r}(|S|)\geq1>0=\vp{r}(x)$.
(${}^2G_2(3)$ is not simple; ${}^2G_2(3)'\cong\PSL(2,8)$ is covered by
Theorem~\ref{thm:psl2}.)
\end{proof}

\subsection{Untwisted types with graph symmetry}

\begin{theorem}\label{thm:graph}
Let $S$ be one of
\[
\PSL(n,q)\ (n\geq3),\quad P\Omega^+(2n,q)=D_n(q)\ (n\geq4),\quad
E_6(q),
\]
\[
\Sp(4,q)\ (q=2^f,\ f\geq2),\quad F_4(q)\ (q=2^f),\quad
G_2(q)\ (q=3^f).
\]
Then $S$ is excluded.
\end{theorem}

\begin{proof}
We keep the standard-parabolic notation $Q_J$ of Lemma~\ref{lem:P}. The
\emph{graph part} of $X$ is the image of $X$ in the quotient of
$\Out(S)$ by its inner-diagonal and field part: a subgroup of
$C_2=\langle\delta\rangle$ for $\PSL(n,q)$ ($n\geq3$), $D_n$ ($n\geq5$),
$E_6$; a subgroup of $S_3$ (triality) for $D_4$; while for $\Sp(4,2^f)$,
$F_4(2^f)$, $G_2(3^f)$, $\Out(S)$ is cyclic of order $2f$ generated by
the exceptional graph-field automorphism $\rho$ with $\rho^2=\varphi_p$
\cite[Def.~3.1(b), p.~31, and Thm.~3.4, p.~32]{BrotoMollerOliver2019},
and ``$X$ has graph part'' means the image of
$X$ in $\Out(S)$ is not contained in $\langle\rho^2\rangle$. An
automorphism acts on the building permuting vertex types by a diagram
symmetry; automorphisms with trivial graph part preserve every type. By
Convention~\ref{conv:X} we may normalize the graph part up to
$\Aut(S)$-conjugacy; for $D_4$ this lets us assume a $C_2$ graph part is
generated by the symmetry fixing nodes $1,2$ and swapping $3,4$.

\emph{Case A: $X$ with trivial graph part.} Every parabolic class is
$X$-stable, and the argument of Theorem~\ref{thm:nograph} applies. The selected
pairs and obstruction primes are given below; the $E_6(q)$ entry refers forward
to a construction valid for every $X$.
\begin{itemize}
\item \emph{$\PSL(n,q)$: pair $([P_1],[P_2])$, $r$ a primitive prime divisor
of $p^{nf}-1$.} Then $r\mid\Phi_n(q)\mid q^n-1$, and this cyclotomic
factor divides $|S|\cdot d$, where $d=\gcd(n,q-1)$. The Levi factors,
of $\GL_{n-1}(q)$-type and $\GL_2{\times}\GL_{n-2}$-type, have
cyclotomic exponents at most $(n-1)f<nf$. Moreover, $x\mid2df$. The
congruence $r\equiv1\pmod{nf}$ rules out $r\mid2f$, since
$r>nf\geq3f>2f$, and also rules out $r\mid d$, since that would imply
$r\mid q-1$ and hence $\ord_r(p)\mid f$. The Zsigmondy exceptions occur
when $p^{nf}=2^6$. If $(n,f)=(3,2)$, then $S=\PSL(3,4)$, which is in the
machine base for all ten $X$. If $(n,f)=(6,1)$, then $S=\PSL(6,2)$;
take instead $r=31$ ($\ord_{31}(2)=5$) and the pair
$([P_2],[P_3])$. Their Levi types $\GL_2{\times}\GL_4$ and
$\GL_3{\times}\GL_3$ have exponents at most $4<5$, so
$31\nmid|P_2||P_3|$, whereas $\vp{31}(|\PSL(6,2)|)=1$ and $x\mid2$.

\item \emph{$D_n(q)$, $n\geq4$: pair $([P_1],[P_2])$ (singular-point and
singular-line stabilizers).} Let $r$ be a primitive prime divisor of
$p^{2(n-1)f}-1$. Then $r\mid\Phi_{2n-2}(q)$, which divides
$q^{n(n-1)}(q^n-1)\prod_{i=1}^{n-1}(q^{2i}-1)$. The Levi factors, of
$D_{n-1}$-type and $A_1{\times}D_{n-2}$-type, have exponents at most
$\max((n-1)f,2(n-2)f)<2(n-1)f$. Moreover, $x\mid24f$, accounting for
diagonal, field, and graph factors of orders at most $4$, $f$, and $6$,
respectively. Since $r\equiv1\pmod{2(n-1)f}$ and $2(n-1)\geq6$, we
have $r>6f$ and $r\nmid24$. The only Zsigmondy exception is
$p^{2(n-1)f}=2^6$, which gives $(n,f)=(4,1)$ and
$S=\Omega^+(8,2)$; Case~B3 handles this group for every $X$.

\item \emph{$E_6(q)$.} The pair used in Case~B4 is stable for every $X$.

\item \emph{$\Sp(4,2^f)$.} Take the pair $([P_1],[P_2])$, and let $r$
be a primitive prime divisor of $2^{4f}-1$. Such a prime exists because
$4f\geq8$. It divides
$\Phi_4(q)=q^2+1$. The Levi factors $\Sp_2(q)$ and $\GL_2(q)$ have
exponents at most $2f<4f$; here $x\mid f$ and $r>4f$.

\item \emph{$F_4(2^f)$.} Take the pair $([P_1],[P_2])$, and let $r$ be a
primitive prime divisor of $2^{12f}-1$. It divides $\Phi_{12}(q)$. All
proper Levi factors, of types
$B_3$, $C_3$, $A_2{\times}A_1$, and $A_1{\times}A_2$, have exponents
at most $6f<12f$, and $x\mid f<r$. This completes the $p=2$ case
excluded from Theorem~\ref{thm:nograph}.

\item \emph{$G_2(3^f)$.} Take the pair $([P_1],[P_2])$, and let $r$ be a
primitive prime divisor of $3^{6f}-1$. It divides $\Phi_6(q)$. Levi factors
of type $A_1$ have
exponents at most $2f<6f$, and $x\mid f<r$. This is the $p=3$ case of
$G_2$.
\end{itemize}

\emph{Case B: $X$ with nontrivial graph part.} Novelty classes below are
standard parabolics $Q_J$ handled by Lemma~\ref{lem:P}: each $[Q_J]$ is
$X$-stable because $J$ is invariant under the relevant diagram symmetry
(an automorphism with graph part $\gamma$ maps $[Q_J]$ to
$[Q_{\gamma(J)}]$), and its proper parabolic overgroups fall into
nontrivial $\gamma$-orbits, giving hypothesis (b) of Lemma~\ref{lem:P};
so $B_{Q_J}$ is maximal in $G$ by Lemma~\ref{lem:B}, and pairs of
distinct classes are non-conjugate by Lemma~\ref{lem:C}. The primes $r$
are those of Case~A unless stated; $r$-avoidance for $Q_J$ follows from
its own Levi order (a product of cyclotomic factors of the Levi of type
$J$ and powers of $q-1$), as listed.

\emph{B1. $\PSL(3,q)$, $q\geq5$.} (For $q\leq4$ the group is in the
machine base, including all ten $X$ for $\PSL(3,4)$.) The pair
consists of $[B]$ and $[N_T]$, where
$N_T:=N_S(T)\cong((q-1)^2/d).S_3$ is the normalizer of a maximal split
torus and $d=\gcd(3,q-1)$. The Borel class is a novelty: its proper
overgroups $P_1,P_2$ are swapped by $\delta$, which exchanges the
points and lines of $PG(2,q)$, so Lemma~\ref{lem:P} applies. The
torus normalizer is maximal for $q\geq5$
\cite{Mitchell1911,Hartley1925}, \cite[Table~8.3, p.~378]{BHRD} (class
$\mathcal{C}_2$). Its conjugates form a single class preserved by every
automorphism. Let $r$ be a primitive prime divisor of $p^{3f}-1$. Then
$r\mid\Phi_3(q)=q^2+q+1$ and $\vp{r}(|S|)\geq1$. Primitivity gives
$r\nmid|B|=q^3(q-1)^2/d$ and $r\nmid q-1$, and
$r\nmid|N_T|=6(q-1)^2/d$ follows because $r\equiv1\pmod3$ forces
$r\geq7$, so $r\nmid6$. Finally, $x\mid2df$: $r>3f$ rules out
$r\mid2f$, and $r\geq7$ rules out $r\mid d=3$. The Zsigmondy
exception $p^{3f}=2^6$ is $\PSL(3,4)$, in the machine base.

\emph{B2. $\PSL(n,q)$, $n\geq4$.} For $n=4$ the pair is
$([P_2],[Q_{\{2\}}])$.
Node $2$ is $\delta$-fixed, so $[P_2]$ is an $X$-stable \emph{maximal}
class (Levi $A_1{\times}A_1$, exponents $\leq2f$), while $Q_{\{2\}}$,
the stabilizer of an incident point--hyperplane flag, is a novelty
whose proper overgroups
$P_1,P_3$ are swapped by $\delta$ (Lemma~\ref{lem:P}); its Levi ($A_1$
plus torus) has exponents $\leq2f$. Let $r$ be a primitive prime divisor
of $p^{4f}-1$.
It avoids both orders because $2f<4f$, and $\vp{r}(|S|)\geq1$. Here
$x\mid2df$, where $d=\gcd(4,q-1)$; the congruence
$r\equiv1\pmod{4f}$ excludes both $r\mid2f$ and $r\mid d$. There is no
Zsigmondy exception ($2^{4f}=2^6$ is
impossible).
For $n\geq5$ the pair consists of the two novelties
$([Q_{J_1}],[Q_{J_2}])$, where
$J_1:=\Delta\setminus\{1,n{-}1\}$ and
$J_2:=\Delta\setminus\{2,n{-}2\}$. The subgroups $Q_{J_1}$ and $Q_{J_2}$ are,
respectively, the stabilizers of an incident point--hyperplane flag and an
incident line--$(n{-}2)$-space flag. Both deleted node sets $\{1,n{-}1\}$ and
$\{2,n{-}2\}$
are $\delta$-orbits, so both $[Q_{J_i}]$ are $X$-stable; the proper
overgroups of $Q_{J_1}$ are $P_1,P_{n-1}$, swapped by $\delta$, and
those of
$Q_{J_2}$ are $P_2,P_{n-2}$, swapped by $\delta$ and distinct classes
since $n\geq5$.
The Levi types are the blocks $(1,n{-}2,1)$, i.e.\ $\GL_{n-2}$ plus
torus, with
exponents $\leq(n{-}2)f$, and the blocks $(2,n{-}4,2)$, with exponents
$\leq\max(2,n{-}4)f$. The prime $r$ of Case~A avoids both, and
$\vp{r}(x)=0$ as there. The Zsigmondy exception is $(n,f)=(6,1)$,
$S=\PSL(6,2)$: take $r=31$ again; the two Levi orders involve only
$\GL_4(2)$- and $\GL_2(2)$-factors and $2$-powers, so $31$ divides neither,
while $\vp{31}(|S|)=1>0=\vp{31}(2)$.

\emph{B3. $D_n(q)$.} For $n\geq5$ the graph part is at most the $C_2$
swapping the two spinor nodes $n{-}1,n$ and fixing $1,\dots,n{-}2$: the
Case~A pair $([P_1],[P_2])$ is $X$-stable for \emph{every} $X$ and
Case~A's argument goes through unchanged (its $x$-bound already allowed
the graph factor). For $n=4$, if the graph part is trivial or the
normalized $C_2$ (fixing nodes $1,2$), the pair $([P_1],[P_2])$ is
$X$-stable and Case~A applies. If the graph part contains a triality,
take the pair $([P_2],[Q_{\{2\}}])$: node $2$ is fixed by all of $S_3$,
so $[P_2]$ is an $X$-stable maximal class (Levi
$A_1{\times}A_1{\times}A_1$, exponents $\leq2f$), and $Q_{\{2\}}$ (Levi
$A_1$ plus torus, exponents $\leq2f$) has proper overgroups
$Q_{\{1,2\}},Q_{\{2,3\}},Q_{\{2,4\}}$ and $P_4,P_3,P_1$, which fall
into the two regular $\langle\text{triality}\rangle$-orbits induced by
the permutation of the end nodes $\{1,3,4\}$; none is $X$-stable
(Lemma~\ref{lem:P}). Let $r$ be a primitive prime divisor of $p^{6f}-1$.
It avoids both classes, $\vp{r}(|S|)\geq1$, and $r\nmid x$ as in Case~A.
\emph{Zsigmondy exception ($f=1$, $p=2$): $S=\Omega^+(8,2)$.} Here
$x\mid24$; take $r=5=\Phi_4(2)$, so
$\vp{5}(|S|)=\vp{5}\bigl(2^{12}(2^6-1)(2^4-1)^2(2^2-1)\bigr)=2$ and
$\vp{5}(x)\leq\vp{5}(24)=0$. In the triality case the pair
$([P_2],[Q_{\{2\}}])$ still serves: both
Levi orders are $\{2,3\}$-numbers, so $\eta=2-0-0=2>0$. In the
trivial/$C_2$ case the pair $([P_1],[P_2])$ gives
$\vp{5}(|P_1|)=\vp{5}(|\SL_4(2)|)=1$ (Levi
$D_3\cong A_3$) and $\vp{5}(|P_2|)=0$, so $\eta=2-1-0=1>0$.

\emph{B4. $E_6(q)$, all $X$.} The diagram involution fixes nodes $2$
and $4$ (Bourbaki numbering: $\delta$ swaps $1\leftrightarrow6$,
$3\leftrightarrow5$). The pair is $([P_2],[P_4])$: two $X$-stable
classes of \emph{maximal} parabolics for every $X$. Let $r$ be a primitive
prime divisor of $p^{12f}-1$. Then $r\mid\Phi_{12}(q)$, which divides
$|S|\cdot d$. For these classes,
the Levi of $P_2$ has type $A_5$ and exponents at most $6f$; that of $P_4$ has type
$A_2{\times}A_1{\times}A_2$ and exponents at most $3f$. Both bounds are
less than $12f$. Finally, $x\mid2df\leq6f<r$, where
$d=\gcd(3,q-1)$. There is no Zsigmondy exception. This simultaneously covers
Case~A for $E_6$.

\emph{B5. $\Sp(4,q)$, $q=2^f$, $f\geq2$, $X$ with graph part.} (For
$f=1$ the group $\Sp(4,2)$ is not simple; its derived group $A_6$ is
in the machine base.) The pair consists of $[B]$ and $[V]$. The Borel
class is a novelty: its proper overgroups $P_1,P_2$, the point
and totally-isotropic-line stabilizers, are swapped by $\rho$, so
Lemma~\ref{lem:P} applies. For the second class, put
$V:=\Fix(\tilde\rho^{\,f})\cap S$, where $\tilde\rho$ is the Steinberg
endomorphism of the ambient algebraic group with $\tilde\rho^2=\varphi_2$
and $\tilde\rho^{\,2f}=F_q$. Thus $V=\Sz(q)$ for odd $f$,
whereas for even $f$ the group $V=\Sp(4,2^{f/2})$ is the subfield
subgroup over $\F_{2^{f/2}}$ corresponding to the degree-two extension
$\F_{2^f}/\F_{2^{f/2}}$. In both cases $V$ is maximal
\cite[Table~8.14, p.~384]{BHRD}, \cite{Suzuki1962}.

\emph{Stability of $[V]$.}
Here $\Aut(S)=\Inn(S)\langle\rho\rangle$ with $\rho=\tilde\rho|_S$, and
$\tilde\rho$ commutes with its own power $\tilde\rho^{\,f}$, so
$\rho(V)=V$; inner automorphisms preserve $[V]$ trivially, so
$[V]$ is $\Aut$-stable.

\emph{Primes.} \emph{Case $f$ odd, $2f\neq6$.} Let $r$ be a primitive
prime divisor of $2^{2f}-1$. Such a prime divides $q+1$, so
$\vp{r}(|S|)=\vp{r}\bigl((q^2-1)(q^4-1)\bigr)=2\vp{r}(q+1)\geq2$,
while $\vp{r}(|B|)=\vp{r}(q^4(q-1)^2)=0$ and $\vp{r}(|\Sz(q)|)=0$ since
$\gcd(q+1,\,q(q^2+1)(q-1))=1$ for $q$ even; finally $x\mid2f<r$.
If $2f=6$ (so $q=8$), take $r=3$; then
$\vp{3}(|\Sp(4,8)|)=\vp{3}(63\cdot4095)=4>1\geq\vp{3}(x)$, with
$\vp{3}(|B|)=\vp{3}(|\Sz(8)|)=0$.

\emph{Case $f$ even.} Let $r$ be a primitive prime divisor of
$2^{4f}-1$. Such a prime exists because $4f\geq8$; it divides $q^2+1$,
so $\vp{r}(|S|)\geq1$;
$|\Sp(4,q_0)|=q_0^4(q_0^2-1)(q_0^4-1)$ with $q_0^2=q$ has cyclotomic
exponents $\leq2f<4f$, so $\vp{r}(|V|)=0$ and $\vp{r}(|B|)=0$; again
$x\mid2f<r$. The sweep-K3 certificate for $\Sp(4,4)$ records the same
construction: the pair (Borel, $\Sp(4,2)$) at $r=5=\Phi_4(2)$ and
$r=17=\Phi_8(2)$.

\emph{B6. $F_4(q)$, $q=2^f$, $X$ with graph part.} Here $\rho$ reverses the
diagram $1\!-\!2\Rightarrow3\!-\!4$, swapping $1\leftrightarrow4$
and $2\leftrightarrow3$. The pair consists of the two novelties
$([Q_{J_1}],[Q_{J_2}])$, where
$J_1:=\{2,3\}$ and $J_2:=\{1,4\}$, both $\rho$-invariant node sets.
The proper overgroups of $Q_{J_1}$ are $Q_{\{1,2,3\}}=P_4$ and
$Q_{\{2,3,4\}}=P_1$, swapped by $\rho$; those of $Q_{J_2}$ are
$Q_{\{1,2,4\}}$
and $Q_{\{1,3,4\}}$, also swapped by $\rho$ and of distinct types,
hence in distinct
classes. So Lemma~\ref{lem:P} applies to both. The corresponding Levi types
are $B_2$ and $A_1{\times}A_1$. The first is isomorphic to $\Sp(4,q)$ and
has exponents at most $4f$, while the second has exponents at most $2f$.
Let $r$ be a primitive prime divisor of $2^{12f}-1$. It avoids both, and
$\vp{r}(|S|)\geq1$ while $x\mid2f<12f<r$. There are no exceptions.

\emph{B7. $G_2(q)$, $q=3^f$, $X$ with graph part.} Here $\rho$ swaps the two
nodes (long and short roots) and $\rho^2=\varphi_3$. The pair again
consists of $[B]$, a novelty whose
overgroups $P_1,P_2$ are swapped by $\rho$ (Lemma~\ref{lem:P}), and
$[V]$, where
$V:=\Fix(\tilde\rho^{\,f})\cap S$ as in B5. Thus $V={}^2G_2(q)$ for
$f$ odd, maximal in $G_2(q)$ \cite[Thm.~A, p.~33]{Kleidman1988} (for $f=1$,
$V={}^2G_2(3)\cong\mathrm{P}\Gamma\mathrm{L}_2(8)$, maximal in $G_2(3)$ by the
\textsc{Atlas} \cite{ATLAS} and \cite[Thm.~A, p.~33]{Kleidman1988}), and
$V=G_2(3^{f/2})$ for even $f$. In the latter case, $V$ is the subfield
subgroup over $\F_{3^{f/2}}$ corresponding to the degree-two extension
$\F_{3^f}/\F_{3^{f/2}}$, and it too is maximal
\cite[Thm.~A, p.~33]{Kleidman1988}. The class $[V]$ is
$\Aut$-stable by the commuting-endomorphism argument of B5, since
$\Out(G_2(3^f))=\langle\rho\rangle\cong C_{2f}$.

\emph{Primes.} \emph{Case $f$ odd.} Let $r$ be a primitive prime divisor
of $3^{3f}-1$. Such a prime exists for all $f\geq1$, because
$p=3\neq2$
excludes the $2^6$ exception ($f=1$ gives $r=13\mid3^3-1$). Then
$r\mid\Phi_3(q)=q^2+q+1$, which divides $q^6-1$. Since
$|S|=q^6(q^6-1)(q^2-1)$, it follows that $r\mid|S|$. Primitivity gives
$r\nmid|B|=q^6(q-1)^2$ and
$r\nmid q-1$ (as $r\mid 3^f-1$ would force $\ord_r(3)\leq f<3f$),
while $r\nmid|{}^2G_2(q)|=q^3(q^3+1)(q-1)$ because $r\mid q^3-1$ and
$\gcd(q^3-1,q^3+1)=2<r$. Since $x\mid2f$ and
$r\equiv1\pmod{3f}$, we have $r\geq3f+1>2f$; hence
$\vp{r}(x)=0<1\leq\vp{r}(|S|)$.

\emph{Case $f$ even.} Let $r$ be a primitive prime divisor of
$3^{6f}-1$. Such a prime divides $\Phi_6(q)=q^2-q+1$, which divides
$q^3+1$; hence
$r\nmid|B|$ and $r\nmid|G_2(3^{f/2})|=q^3(q^3-1)(q-1)$, while
$x\mid2f<6f<r$. ($G_2(3)$ additionally carries a full machine
certificate for both $X$.)

In every case Theorem~\ref{thm:D} excludes $S^k$ for every $k\geq2$.
\end{proof}

\section{Proof of the Main Theorem}\label{sec:proof}

\begin{proposition}[Coverage]\label{prop:coverage}
Every non-abelian finite simple group is excluded in the sense of
Convention~\ref{conv:excluded}.
\end{proposition}

\begin{proof}
By the classification of finite simple groups
\cite[Ch.~1, Table~I, pp.~8--10]{GLS1}, $S$ is alternating,
of Lie type, or one of the $26$ sporadic groups. The cases are as follows.
The simplicity exceptions are
those in \cite[Thm.~2.2.7(a)]{GLS3}; the low-rank classical
identifications and the exceptional derived-group identifications are
recorded in \cite[\S2.4, pp.~xi--xii, and \S3.5, p.~xv]{ATLAS}
(with the individual \textsc{Atlas} entries pinpointed in the source map in
\cite{Repo}).
\begin{itemize}
\item $A_n$ ($n\geq5$): Theorem~\ref{thm:an}.
\item $\PSL(2,q)$ ($q\geq4$): Theorem~\ref{thm:psl2}.
\item $\PSL(n,q)$ ($n\geq3$), $D_n(q)$ ($n\geq4$), $E_6(q)$:
Theorem~\ref{thm:graph}.
\item $\PSp(2n,q)$: $n\geq3$, or $n=2$ with $q$ odd:
Theorem~\ref{thm:nograph}; $\Sp(4,2^f)$, $f\geq2$:
Theorem~\ref{thm:graph} ($\Sp(4,2)$ is not simple).
\item $\Omega(2n{+}1,q)$, $q$ odd, $n\geq3$: Theorem~\ref{thm:nograph}
(for $q$ even, $B_n(q)\cong C_n(q)$, already listed;
$\Omega(5,q)\cong\PSp(4,q)$).
\item $E_7(q)$, $E_8(q)$; $F_4(q)$, $p\neq2$; $G_2(q)$, $p\neq3$:
Theorem~\ref{thm:nograph}. $F_4(2^f)$ and $G_2(3^f)$:
Theorem~\ref{thm:graph}. ($G_2(2)$ and ${}^2G_2(3)$ are not simple;
their derived groups $\PSU(3,3)$ and $\PSL(2,8)$ are covered.)
\item $\PSU(n,q)$: $n=3$: Theorem~\ref{thm:twisted1}; $n\geq4$:
Theorem~\ref{thm:twisted2}.
\item ${}^2D_n(q)$ ($n\geq4$), ${}^3D_4(q)$, ${}^2E_6(q)$,
${}^2F_4(2^f)$ ($f$ odd $\geq3$): all four families are covered by
Theorem~\ref{thm:twisted2}; the Tits group ${}^2F_4(2)'$ by
Proposition~\ref{prop:sporadic}. (${}^2D_2$, ${}^2D_3$ are
$\PSL(2,q^2)$, $\PSU(4,q)$, already listed.)
\item $\Sz(2^f)$ ($f$ odd $\geq3$), ${}^2G_2(3^f)$ ($f$ odd $\geq3$):
Theorem~\ref{thm:twisted1}.
\item Sporadic groups: Proposition~\ref{prop:sporadic}.
\end{itemize}
Every finite simple group appears (repetitions from exceptional
isomorphisms are harmless), so the proposition follows.
\end{proof}

\begin{proof}[Proof of Theorem~\ref{thm:main}]
Suppose not, and let $G$ have least order among the non-soluble groups
with property $\mathrm{P}$. By Proposition~\ref{prop:min}, $G$ has a
unique minimal normal subgroup $N=S^k$, $S$ non-abelian simple, and
after conjugation $N\leq G\leq X\wr S_k$ as in Convention~\ref{conv:X}.
If $k=1$, then $G$ is almost simple, contrary to the theorem of
Tikhonenko--Tyutyanov \cite{TT2010}. Hence $k\geq2$. By
Proposition~\ref{prop:coverage} there are classes
$[V]\neq[W]$, maximal in $\Pos_X(S)$ and normally saturating, and a
prime $p$ with $\vp{p}(|S|)-\vp{p}(|V|)-\vp{p}(|W|)>\vp{p}(x)$. By
Theorem~\ref{thm:D}, $G$ does not have property $\mathrm{P}$, a
contradiction.
\end{proof}

\begin{corollary}[Supplementary reproof]
The present criterion also proves that no almost simple group has property
$\mathrm{P}$.
\end{corollary}

\begin{proof}
Apply Theorem~\ref{thm:Dprime} with the same witnesses used in
Proposition~\ref{prop:coverage}.
\end{proof}

\begin{remark}
The argument proves the following stronger statement: for \emph{every}
finite group $G$ with a unique minimal
normal subgroup $S^k$ ($S$ non-abelian simple, $k\geq2$), some pair of
non-conjugate maximal subgroups of $G$ fails to factorize $G$. The
reduction of \S\ref{sec:reduction} is needed only to pass from an
arbitrary non-soluble group to this configuration.
\end{remark}

\subsection*{Data Availability}
All verification code, including the GAP scripts, and all certificate
logs required to reproduce the
machine-verified claims of \S\ref{sec:machine}, together with the
arithmetic verification records for the family proofs of
\S\ref{sec:families}, are
archived in the versioned release \cite{Repo}. The release includes
SHA-256 checksums of every certificate log, a \path{CITATION.cff}
file, a single command that runs the quick verification suite, and
separate documented commands for full GAP certificate regeneration and
byte comparison, Lean and Rocq builds, and mutation-test
reproduction. Its README maps each claim of this paper to the script
and certificate that verify it. The priority-search record and source
ledger supporting the discussion in the Introduction are also included.

\subsection*{Funding}
This research received no external funding.

\subsection*{Competing interests}
The author declares no competing interests.

\subsection*{Declaration of generative AI and AI-assisted technologies
in the research and writing process}
During this project the author used Anthropic's Claude (2026
Claude~5-family models) and OpenAI Codex for mathematical exploration,
generation and critical evaluation of candidate arguments,
formal-proof development,
drafting and debugging the GAP and Python code, literature-search assistance,
verification-workflow construction, and manuscript organization and editing. The
author reviewed the mathematical arguments, reran the verification
scripts or checked their outputs against the committed certificate
logs, verified the cited references, and edited the final manuscript. All
machine-verified claims are backed by the committed, independently
checkable certificates described in \S\ref{sec:machine}. The author
takes full responsibility for the content of this article.

\begin{thebibliography}{99}

\bibitem{BHRD}
J.~N.~Bray, D.~F.~Holt, C.~M.~Roney-Dougal,
\emph{The Maximal Subgroups of the Low-Dimensional Finite Classical
Groups}, London Math.\ Soc.\ Lecture Note Ser.\ 407, Cambridge Univ.\
Press, 2013; DOI 10.1017/CBO9781139192576.

\bibitem{CTblLib}
T.~Breuer, \emph{The GAP Character Table Library}, GAP package,
version 1.3.11 (2025),
\path{https://www.math.rwth-aachen.de/~Thomas.Breuer/ctbllib}.

\bibitem{BreuerMalleOBrien2017}
T.~Breuer, G.~Malle, E.~A.~O'Brien, Reliability and reproducibility of
Atlas information, in \emph{Finite Simple Groups: Thirty Years of the
Atlas and Beyond}, Contemp.\ Math.\ \textbf{694} (2017), 21--32;
DOI 10.1090/conm/694/13960.

\bibitem{BrotoMollerOliver2019}
C.~Broto, J.~M.~M\o ller, B.~Oliver, \emph{Automorphisms of Fusion
Systems of Finite Simple Groups of Lie Type}, Mem.\ Amer.\ Math.\ Soc.\
\textbf{262}, no.~1267 (2019); DOI 10.1090/memo/1267.

\bibitem{CameronNotes}
P.~J.~Cameron, \emph{Notes on Classical Groups}, Queen Mary and
Westfield College M.Sc.\ course notes, 2000,
\path{https://webspace.maths.qmul.ac.uk/p.j.cameron/class_gps/cg.pdf}.

\bibitem{Carter}
R.~W.~Carter, \emph{Simple Groups of Lie Type}, Wiley, London, 1972.

\bibitem{Carter1965}
R.~W.~Carter, Simple groups and simple Lie algebras,
\emph{J.\ London Math.\ Soc.} \textbf{40} (1965), 193--240;
DOI 10.1112/jlms/s1-40.1.193.

\bibitem{ATLAS}
J.~H.~Conway, R.~T.~Curtis, S.~P.~Norton, R.~A.~Parker, R.~A.~Wilson,
\emph{Atlas of Finite Groups}, Oxford Univ.\ Press, 1985.

\bibitem{Dickson}
L.~E.~Dickson, \emph{Linear Groups with an Exposition of the Galois
Field Theory}, Teubner, Leipzig, 1901.

\bibitem{DietrichEtAl2026}
H.~Dietrich, M.~Lee, A.~Pisani, T.~Popiel, Explicit construction of the
maximal subgroups of the Monster, \emph{J.\ Algebra} \textbf{689} (2026),
862--895; DOI 10.1016/j.jalgebra.2025.10.016.

\bibitem{GAP}
The GAP Group, \emph{GAP -- Groups, Algorithms, and Programming},
Version 4.16.0, 2026, \path{https://www.gap-system.org}.

\bibitem{GLS1}
D.~Gorenstein, R.~Lyons, R.~Solomon, \emph{The Classification of the
Finite Simple Groups, Number 1}, Math.\ Surveys Monogr.\ 40.1, Amer.\
Math.\ Soc., 1994.

\bibitem{GLS3}
D.~Gorenstein, R.~Lyons, R.~Solomon, \emph{The Classification of the
Finite Simple Groups, Number 3}, Math.\ Surveys Monogr.\ 40.3, Amer.\
Math.\ Soc., 1998.

\bibitem{Hartley1925}
R.~W.~Hartley, Determination of the ternary collineation groups whose
coefficients lie in the $GF(2^n)$, \emph{Ann.\ of Math.} \textbf{27}
(1925), 140--158.

\bibitem{Hiss2011}
G.~Hiss, Finite groups of Lie type and their representations, in
\emph{Groups St Andrews 2009 in Bath}, vol.~1, Cambridge Univ.\ Press,
2011, 1--40; DOI 10.1017/CBO9780511842467.003.

\bibitem{Huppert}
B.~Huppert, \emph{Endliche Gruppen I}, Grundlehren Math.\ Wiss.\ 134,
Springer, Berlin, 1967.

\bibitem{Jones2007}
L.~Jones, Variations on a theme of Sierpi\'nski,
\emph{J.\ Integer Seq.} \textbf{10} (2007), Article 07.4.4, 15 pp.

\bibitem{Kourovka}
E.~I.~Khukhro, V.~D.~Mazurov (eds.), \emph{The Kourovka Notebook:
Unsolved Problems in Group Theory}, 21st ed., Sobolev Inst.\ Math.,
Novosibirsk, 2026; July 2026 revision, p.~39, Problem 10.34
(V.~S.~Monakhov), \path{https://kourovkanotebookorg.wordpress.com/}.

\bibitem{Kleidman1988}
P.~B.~Kleidman, The maximal subgroups of the Chevalley groups $G_2(q)$
with $q$ odd, the Ree groups ${}^2G_2(q)$, and their automorphism
groups, \emph{J.\ Algebra} \textbf{117} (1988), 30--71;
DOI 10.1016/0021-8693(88)90239-6.

\bibitem{LPS1987}
M.~W.~Liebeck, C.~E.~Praeger, J.~Saxl, A classification of the maximal
subgroups of the finite alternating and symmetric groups,
\emph{J.\ Algebra} \textbf{111} (1987), 365--383;
DOI 10.1016/0021-8693(87)90223-7.

\bibitem{LPS1990}
M.~W.~Liebeck, C.~E.~Praeger, J.~Saxl, \emph{The Maximal Factorizations
of the Finite Simple Groups and Their Automorphism Groups}, Mem.\ Amer.\
Math.\ Soc.\ \textbf{86}, no.~432 (1990), 151 pp.

\bibitem{LucchiniMorini2002}
A.~Lucchini, F.~Morini, On the probability of generating finite groups
with a unique minimal normal subgroup, \emph{Pacific J.\ Math.}
\textbf{203} (2002), 429--440; DOI 10.2140/pjm.2002.203.429.

\bibitem{Malle1991}
G.~Malle, The maximal subgroups of ${}^2F_4(q^2)$,
\emph{J.\ Algebra} \textbf{139} (1991), 52--69;
DOI 10.1016/0021-8693(91)90283-E.

\bibitem{Mitchell1911}
H.~H.~Mitchell, Determination of the ordinary and modular ternary
linear groups, \emph{Trans.\ Amer.\ Math.\ Soc.} \textbf{12} (1911),
207--242.

\bibitem{Montinaro2024}
A.~Montinaro, On the symmetric $2$-$(v,k,\lambda)$ designs with a
flag-transitive point-imprimitive automorphism group,
\emph{J.\ Algebra} \textbf{653} (2024), 54--101;
DOI 10.1016/j.jalgebra.2024.05.002.

\bibitem{Nagura}
J.~Nagura, On the interval containing at least one prime number,
\emph{Proc.\ Japan Acad.} \textbf{28} (1952), 177--181;
DOI 10.3792/pja/1195570997.

\bibitem{RoneyDougal2021}
C.~M.~Roney-Dougal, Maximal subgroups of finite simple groups:
classifications and applications, in \emph{Surveys in Combinatorics
2021}, London Math.\ Soc.\ Lecture Note Ser.\ 470, Cambridge Univ.\ Press,
2021, 343--370; DOI 10.1017/9781009036214.010.

\bibitem{Repo}
R.~Sater, \emph{Kourovka 10.34: proofs, GAP scripts and machine
certificates}, version 1.0.8, Zenodo, 2026;
concept DOI 10.5281/zenodo.21709124 (resolves to the latest version);
\path{https://github.com/RichieSater/kourovka-10-34/releases/tag/v1.0.8}.

\bibitem{Steinberg}
R.~Steinberg, Automorphisms of finite linear groups,
\emph{Canad.\ J.\ Math.} \textbf{12} (1960), 606--615.

\bibitem{Suzuki1962}
M.~Suzuki, On a class of doubly transitive groups,
\emph{Ann.\ of Math.} \textbf{75} (1962), 105--145;
DOI 10.2307/1970423.

\bibitem{Tchakerian1986}
K.~B.~Tchakerian, The maximal subgroups of the Tits simple group,
\emph{Pliska Stud.\ Math.\ Bulgar.} \textbf{8} (1986), 85--93.

\bibitem{TT2010}
T.~V.~Tikhonenko, V.~N.~Tyutyanov, On factorization of finite almost
simple groups by nonconjugate maximal subgroups, \emph{Sib.\ Math.\ J.}
\textbf{51}, no.~1 (2010), 174--177; Russian original pp.~212--216;
DOI 10.1007/s11202-010-0018-3; \path{https://www.mathnet.ru/eng/smj2078}.

\bibitem{Wilson1984}
R.~A.~Wilson, The geometry and maximal subgroups of the simple groups
of A.~Rudvalis and J.~Tits, \emph{Proc.\ London Math.\ Soc.}
\textbf{48} (1984), 533--563; DOI 10.1112/plms/s3-48.3.533.

\bibitem{Wilson2009}
R.~A.~Wilson, \emph{The Finite Simple Groups}, Grad.\ Texts in Math.\ 251,
Springer, London, 2009; DOI 10.1007/978-1-84800-988-2.

\bibitem{Wilson2017}
R.~A.~Wilson, Maximal subgroups of sporadic groups, in \emph{Finite
Simple Groups: Thirty Years of the Atlas and Beyond}, Contemp.\ Math.\
\textbf{694} (2017), 57--72; DOI 10.1090/conm/694/13974.

\bibitem{Zenkov1997}
V.~I.~Zenkov, Factorization of finite groups, in \emph{Abstracts of the
International Algebra Conference in Memory of D.~K.~Faddeev},
St.~Petersburg, 1997, p.~202 (in Russian); bibliographic data and the
announcement are recorded in the Russian original of
Tikhonenko--Tyutyanov (2010), p.~212.

\bibitem{ZhangShi2009}
L.~Zhang, W.~Shi, Noncommuting graph characterization of some simple
groups with connected prime graphs, \emph{Int.\ Electron.\ J.\ Algebra}
\textbf{5} (2009), 169--181.

\bibitem{Zsigmondy}
K.~Zsigmondy, Zur Theorie der Potenzreste, \emph{Monatsh.\ Math.\
Phys.} \textbf{3} (1892), 265--284; DOI 10.1007/BF01692444.

\end{thebibliography}

\end{document}
